% Options for packages loaded elsewhere
% Options for packages loaded elsewhere
\PassOptionsToPackage{unicode}{hyperref}
\PassOptionsToPackage{hyphens}{url}
\PassOptionsToPackage{dvipsnames,svgnames,x11names}{xcolor}
%
\documentclass[
  11pt,
  a4paper,
  oneside,
  open=any]{scrbook}
\usepackage{xcolor}
\usepackage[top=25mm,bottom=28mm,left=28mm,right=22mm,headsep=8mm]{geometry}
\usepackage{amsmath,amssymb}
\setcounter{secnumdepth}{5}
\usepackage{iftex}
\ifPDFTeX
  \usepackage[T1]{fontenc}
  \usepackage[utf8]{inputenc}
  \usepackage{textcomp} % provide euro and other symbols
\else % if luatex or xetex
  \usepackage{unicode-math} % this also loads fontspec
  \defaultfontfeatures{Scale=MatchLowercase}
  \defaultfontfeatures[\rmfamily]{Ligatures=TeX,Scale=1}
\fi
\usepackage{lmodern}
\ifPDFTeX\else
  % xetex/luatex font selection
  \setmainfont[]{Latin Modern Roman}
  \setsansfont[]{Latin Modern Sans}
  \setmonofont[]{Latin Modern Mono}
\fi
% Use upquote if available, for straight quotes in verbatim environments
\IfFileExists{upquote.sty}{\usepackage{upquote}}{}
\IfFileExists{microtype.sty}{% use microtype if available
  \usepackage[]{microtype}
  \UseMicrotypeSet[protrusion]{basicmath} % disable protrusion for tt fonts
}{}
\makeatletter
\@ifundefined{KOMAClassName}{% if non-KOMA class
  \IfFileExists{parskip.sty}{%
    \usepackage{parskip}
  }{% else
    \setlength{\parindent}{0pt}
    \setlength{\parskip}{6pt plus 2pt minus 1pt}}
}{% if KOMA class
  \KOMAoptions{parskip=half}}
\makeatother
% Make \paragraph and \subparagraph free-standing
\makeatletter
\ifx\paragraph\undefined\else
  \let\oldparagraph\paragraph
  \renewcommand{\paragraph}{
    \@ifstar
      \xxxParagraphStar
      \xxxParagraphNoStar
  }
  \newcommand{\xxxParagraphStar}[1]{\oldparagraph*{#1}\mbox{}}
  \newcommand{\xxxParagraphNoStar}[1]{\oldparagraph{#1}\mbox{}}
\fi
\ifx\subparagraph\undefined\else
  \let\oldsubparagraph\subparagraph
  \renewcommand{\subparagraph}{
    \@ifstar
      \xxxSubParagraphStar
      \xxxSubParagraphNoStar
  }
  \newcommand{\xxxSubParagraphStar}[1]{\oldsubparagraph*{#1}\mbox{}}
  \newcommand{\xxxSubParagraphNoStar}[1]{\oldsubparagraph{#1}\mbox{}}
\fi
\makeatother

\usepackage{color}
\usepackage{fancyvrb}
\newcommand{\VerbBar}{|}
\newcommand{\VERB}{\Verb[commandchars=\\\{\}]}
\DefineVerbatimEnvironment{Highlighting}{Verbatim}{commandchars=\\\{\}}
% Add ',fontsize=\small' for more characters per line
\usepackage{framed}
\definecolor{shadecolor}{RGB}{241,243,245}
\newenvironment{Shaded}{\begin{snugshade}}{\end{snugshade}}
\newcommand{\AlertTok}[1]{\textcolor[rgb]{0.68,0.00,0.00}{#1}}
\newcommand{\AnnotationTok}[1]{\textcolor[rgb]{0.37,0.37,0.37}{#1}}
\newcommand{\AttributeTok}[1]{\textcolor[rgb]{0.40,0.45,0.13}{#1}}
\newcommand{\BaseNTok}[1]{\textcolor[rgb]{0.68,0.00,0.00}{#1}}
\newcommand{\BuiltInTok}[1]{\textcolor[rgb]{0.00,0.23,0.31}{#1}}
\newcommand{\CharTok}[1]{\textcolor[rgb]{0.13,0.47,0.30}{#1}}
\newcommand{\CommentTok}[1]{\textcolor[rgb]{0.37,0.37,0.37}{#1}}
\newcommand{\CommentVarTok}[1]{\textcolor[rgb]{0.37,0.37,0.37}{\textit{#1}}}
\newcommand{\ConstantTok}[1]{\textcolor[rgb]{0.56,0.35,0.01}{#1}}
\newcommand{\ControlFlowTok}[1]{\textcolor[rgb]{0.00,0.23,0.31}{\textbf{#1}}}
\newcommand{\DataTypeTok}[1]{\textcolor[rgb]{0.68,0.00,0.00}{#1}}
\newcommand{\DecValTok}[1]{\textcolor[rgb]{0.68,0.00,0.00}{#1}}
\newcommand{\DocumentationTok}[1]{\textcolor[rgb]{0.37,0.37,0.37}{\textit{#1}}}
\newcommand{\ErrorTok}[1]{\textcolor[rgb]{0.68,0.00,0.00}{#1}}
\newcommand{\ExtensionTok}[1]{\textcolor[rgb]{0.00,0.23,0.31}{#1}}
\newcommand{\FloatTok}[1]{\textcolor[rgb]{0.68,0.00,0.00}{#1}}
\newcommand{\FunctionTok}[1]{\textcolor[rgb]{0.28,0.35,0.67}{#1}}
\newcommand{\ImportTok}[1]{\textcolor[rgb]{0.00,0.46,0.62}{#1}}
\newcommand{\InformationTok}[1]{\textcolor[rgb]{0.37,0.37,0.37}{#1}}
\newcommand{\KeywordTok}[1]{\textcolor[rgb]{0.00,0.23,0.31}{\textbf{#1}}}
\newcommand{\NormalTok}[1]{\textcolor[rgb]{0.00,0.23,0.31}{#1}}
\newcommand{\OperatorTok}[1]{\textcolor[rgb]{0.37,0.37,0.37}{#1}}
\newcommand{\OtherTok}[1]{\textcolor[rgb]{0.00,0.23,0.31}{#1}}
\newcommand{\PreprocessorTok}[1]{\textcolor[rgb]{0.68,0.00,0.00}{#1}}
\newcommand{\RegionMarkerTok}[1]{\textcolor[rgb]{0.00,0.23,0.31}{#1}}
\newcommand{\SpecialCharTok}[1]{\textcolor[rgb]{0.37,0.37,0.37}{#1}}
\newcommand{\SpecialStringTok}[1]{\textcolor[rgb]{0.13,0.47,0.30}{#1}}
\newcommand{\StringTok}[1]{\textcolor[rgb]{0.13,0.47,0.30}{#1}}
\newcommand{\VariableTok}[1]{\textcolor[rgb]{0.07,0.07,0.07}{#1}}
\newcommand{\VerbatimStringTok}[1]{\textcolor[rgb]{0.13,0.47,0.30}{#1}}
\newcommand{\WarningTok}[1]{\textcolor[rgb]{0.37,0.37,0.37}{\textit{#1}}}

\usepackage{longtable,booktabs,array}
\usepackage{calc} % for calculating minipage widths
% Correct order of tables after \paragraph or \subparagraph
\usepackage{etoolbox}
\makeatletter
\patchcmd\longtable{\par}{\if@noskipsec\mbox{}\fi\par}{}{}
\makeatother
% Allow footnotes in longtable head/foot
\IfFileExists{footnotehyper.sty}{\usepackage{footnotehyper}}{\usepackage{footnote}}
\makesavenoteenv{longtable}
\usepackage{graphicx}
\makeatletter
\newsavebox\pandoc@box
\newcommand*\pandocbounded[1]{% scales image to fit in text height/width
  \sbox\pandoc@box{#1}%
  \Gscale@div\@tempa{\textheight}{\dimexpr\ht\pandoc@box+\dp\pandoc@box\relax}%
  \Gscale@div\@tempb{\linewidth}{\wd\pandoc@box}%
  \ifdim\@tempb\p@<\@tempa\p@\let\@tempa\@tempb\fi% select the smaller of both
  \ifdim\@tempa\p@<\p@\scalebox{\@tempa}{\usebox\pandoc@box}%
  \else\usebox{\pandoc@box}%
  \fi%
}
% Set default figure placement to htbp
\def\fps@figure{htbp}
\makeatother





\setlength{\emergencystretch}{3em} % prevent overfull lines

\providecommand{\tightlist}{%
  \setlength{\itemsep}{0pt}\setlength{\parskip}{0pt}}



 


% Professional textbook-style refinements for Quarto PDF output
\usepackage{booktabs}
\usepackage{longtable}
\usepackage{array}
\usepackage{caption}

% Improve body text rhythm
\setlength{\parindent}{1.2em}
\setlength{\parskip}{0.3em}

% Cleaner heading hierarchy
\setkomafont{disposition}{\bfseries}
\setkomafont{chapter}{\Huge\bfseries}
\setkomafont{section}{\Large\bfseries}
\setkomafont{subsection}{\large\bfseries}

% Better table and figure captions
\captionsetup{font=small,labelfont=bf}

% Subtle running header/footer for print-like structure
\KOMAoptions{headings=normal}
\pagestyle{headings}
\makeatletter
\@ifpackageloaded{tcolorbox}{}{\usepackage[skins,breakable]{tcolorbox}}
\@ifpackageloaded{fontawesome5}{}{\usepackage{fontawesome5}}
\definecolor{quarto-callout-color}{HTML}{909090}
\definecolor{quarto-callout-note-color}{HTML}{0758E5}
\definecolor{quarto-callout-important-color}{HTML}{CC1914}
\definecolor{quarto-callout-warning-color}{HTML}{EB9113}
\definecolor{quarto-callout-tip-color}{HTML}{00A047}
\definecolor{quarto-callout-caution-color}{HTML}{FC5300}
\definecolor{quarto-callout-color-frame}{HTML}{acacac}
\definecolor{quarto-callout-note-color-frame}{HTML}{4582ec}
\definecolor{quarto-callout-important-color-frame}{HTML}{d9534f}
\definecolor{quarto-callout-warning-color-frame}{HTML}{f0ad4e}
\definecolor{quarto-callout-tip-color-frame}{HTML}{02b875}
\definecolor{quarto-callout-caution-color-frame}{HTML}{fd7e14}
\makeatother
\makeatletter
\@ifpackageloaded{bookmark}{}{\usepackage{bookmark}}
\makeatother
\makeatletter
\@ifpackageloaded{caption}{}{\usepackage{caption}}
\AtBeginDocument{%
\ifdefined\contentsname
  \renewcommand*\contentsname{Table of contents}
\else
  \newcommand\contentsname{Table of contents}
\fi
\ifdefined\listfigurename
  \renewcommand*\listfigurename{List of Figures}
\else
  \newcommand\listfigurename{List of Figures}
\fi
\ifdefined\listtablename
  \renewcommand*\listtablename{List of Tables}
\else
  \newcommand\listtablename{List of Tables}
\fi
\ifdefined\figurename
  \renewcommand*\figurename{Figure}
\else
  \newcommand\figurename{Figure}
\fi
\ifdefined\tablename
  \renewcommand*\tablename{Table}
\else
  \newcommand\tablename{Table}
\fi
}
\@ifpackageloaded{float}{}{\usepackage{float}}
\floatstyle{ruled}
\@ifundefined{c@chapter}{\newfloat{codelisting}{h}{lop}}{\newfloat{codelisting}{h}{lop}[chapter]}
\floatname{codelisting}{Listing}
\newcommand*\listoflistings{\listof{codelisting}{List of Listings}}
\makeatother
\makeatletter
\makeatother
\makeatletter
\@ifpackageloaded{caption}{}{\usepackage{caption}}
\@ifpackageloaded{subcaption}{}{\usepackage{subcaption}}
\makeatother
\usepackage{bookmark}
\IfFileExists{xurl.sty}{\usepackage{xurl}}{} % add URL line breaks if available
\urlstyle{same}
\hypersetup{
  pdftitle={AS91946: Interpret and Apply Mathematical and Statistical Information},
  pdfauthor={Anthony R Davidson},
  colorlinks=true,
  linkcolor={Maroon},
  filecolor={Maroon},
  citecolor={Blue},
  urlcolor={Blue},
  pdfcreator={LaTeX via pandoc}}


\title{AS91946: Interpret and Apply Mathematical and Statistical
Information}
\usepackage{etoolbox}
\makeatletter
\providecommand{\subtitle}[1]{% add subtitle to \maketitle
  \apptocmd{\@title}{\par {\large #1 \par}}{}{}
}
\makeatother
\subtitle{NCEA Level 1 Statistics --- Year 11}
\author{Anthony R Davidson}
\date{July 2026}
\begin{document}
\frontmatter
\maketitle

\renewcommand*\contentsname{Table of contents}
{
\hypersetup{linkcolor=}
\setcounter{tocdepth}{1}
\tableofcontents
}
\listoffigures
\listoftables

\mainmatter
\bookmarksetup{startatroot}

\chapter{Welcome}\label{welcome}

\bookmarksetup{startatroot}

\chapter*{Welcome}\label{welcome-1}
\addcontentsline{toc}{chapter}{Welcome}

\markboth{Welcome}{Welcome}

This resource has been developed to support Year 11 students preparing
for \textbf{NCEA Level 1 Achievement Standard 91946}: \emph{Interpret
and apply mathematical and statistical information in context} (4
credits).

\section{What this standard is about}\label{what-this-standard-is-about}

In this standard you will learn to:

\begin{itemize}
\tightlist
\item
  \textbf{Read and interpret} statistical displays, tables, and graphs
\item
  \textbf{Apply} mathematical and statistical ideas to real-world
  problems
\item
  \textbf{Communicate} findings clearly using statistical language
\item
  \textbf{Evaluate} statistical claims made in media and everyday life
\end{itemize}

\section{How to use this book}\label{how-to-use-this-book}

Each chapter builds on the previous one. You will find:

\begin{itemize}
\tightlist
\item
  Clear explanations of key concepts with New Zealand examples
\item
  Interactive R code blocks showing how statistics is done with real
  data
\item
  Practice questions at Achievement, Merit, and Excellence levels
\item
  Glossary boxes highlighting important vocabulary
\end{itemize}

\begin{tcolorbox}[enhanced jigsaw, bottomtitle=1mm, opacityback=0, titlerule=0mm, leftrule=.75mm, opacitybacktitle=0.6, coltitle=black, toprule=.15mm, colbacktitle=quarto-callout-tip-color!10!white, title=\textcolor{quarto-callout-tip-color}{\faLightbulb}\hspace{0.5em}{Achievement Standard at a Glance}, left=2mm, rightrule=.15mm, toptitle=1mm, colframe=quarto-callout-tip-color-frame, breakable, arc=.35mm, bottomrule=.15mm, colback=white]

\begin{longtable}[]{@{}
  >{\raggedright\arraybackslash}p{(\linewidth - 2\tabcolsep) * \real{0.2500}}
  >{\raggedright\arraybackslash}p{(\linewidth - 2\tabcolsep) * \real{0.7500}}@{}}
\toprule\noalign{}
\begin{minipage}[b]{\linewidth}\raggedright
Grade
\end{minipage} & \begin{minipage}[b]{\linewidth}\raggedright
What you need to do
\end{minipage} \\
\midrule\noalign{}
\endhead
\bottomrule\noalign{}
\endlastfoot
\textbf{Achieved} & Interpret and apply statistical information \\
\textbf{Merit} & Interpret and apply statistical information with
justification \\
\textbf{Excellence} & Interpret and apply statistical information with
statistical insight \\
\end{longtable}

\end{tcolorbox}

\section{Tools we use}\label{tools-we-use}

All data visualisations and calculations in this book are produced using
\textbf{R} and rendered with \textbf{Quarto}. You do not need to be able
to write R code to succeed in this standard --- the code blocks are
provided so your teacher can demonstrate concepts live and so you can
explore data yourself if you choose.

\begin{figure}[H]

{\centering \pandocbounded{\includegraphics[keepaspectratio]{index_files/figure-latex/welcome-plot-1.pdf}}

}

\caption{An example of the kind of visualisations we will create in this
book.}

\end{figure}%

\section{Getting started}\label{getting-started}

Head to Chapter~\ref{sec-intro} to begin your statistics journey!

\bookmarksetup{startatroot}

\chapter{What is Statistics?}\label{what-is-statistics}

\bookmarksetup{startatroot}

\chapter{What is Statistics?}\label{sec-intro}

\section{Why statistics matters}\label{why-statistics-matters}

Every day we are surrounded by data. News headlines cite poll results,
sports commentators compare player averages, health campaigns report
risk percentages, and social media shows us trending topics.
\textbf{Statistics} is the science that helps us make sense of all this
information.

\begin{tcolorbox}[enhanced jigsaw, bottomtitle=1mm, opacityback=0, titlerule=0mm, leftrule=.75mm, opacitybacktitle=0.6, coltitle=black, toprule=.15mm, colbacktitle=quarto-callout-note-color!10!white, title=\textcolor{quarto-callout-note-color}{\faInfo}\hspace{0.5em}{Key definition}, left=2mm, rightrule=.15mm, toptitle=1mm, colframe=quarto-callout-note-color-frame, breakable, arc=.35mm, bottomrule=.15mm, colback=white]

\textbf{Statistics} is the practice of collecting, organising,
summarising, analysing, and interpreting data in order to draw
conclusions and make decisions.

\end{tcolorbox}

\section{The statistical enquiry cycle
(PPDAC)}\label{the-statistical-enquiry-cycle-ppdac}

New Zealand's statistics curriculum uses the \textbf{PPDAC cycle} as a
framework for carrying out a statistical investigation:

\begin{longtable}[]{@{}ll@{}}
\toprule\noalign{}
Stage & What you do \\
\midrule\noalign{}
\endhead
\bottomrule\noalign{}
\endlastfoot
\textbf{P}roblem & Pose a statistical question \\
\textbf{P}lan & Plan how to collect data \\
\textbf{D}ata & Collect and record data \\
\textbf{A}nalysis & Explore and summarise the data \\
\textbf{C}onclusion & Interpret results and answer the question \\
\end{longtable}

\begin{figure}[H]

{\centering \pandocbounded{\includegraphics[keepaspectratio]{chapters/01-intro_files/figure-latex/ppdac-diagram-1.pdf}}

}

\caption{The PPDAC cycle --- a framework for statistical
investigations.}

\end{figure}%

\section{Types of statistical
questions}\label{types-of-statistical-questions}

Not all questions are statistical questions. A statistical question
expects \textbf{variability} in the answers.

\begin{tcolorbox}[enhanced jigsaw, bottomtitle=1mm, opacityback=0, titlerule=0mm, leftrule=.75mm, opacitybacktitle=0.6, coltitle=black, toprule=.15mm, colbacktitle=quarto-callout-tip-color!10!white, title=\textcolor{quarto-callout-tip-color}{\faLightbulb}\hspace{0.5em}{Tip}, left=2mm, rightrule=.15mm, toptitle=1mm, colframe=quarto-callout-tip-color-frame, breakable, arc=.35mm, bottomrule=.15mm, colback=white]

\textbf{Statistical question:} ``How tall are the students in Year
11?''\\
\emph{(Heights will vary --- this is statistical!)}

\textbf{Not a statistical question:} ``How tall is Ms Smith?''\\
\emph{(Only one answer --- not statistical.)}

\end{tcolorbox}

\subsection{Practice identifying statistical
questions}\label{practice-identifying-statistical-questions}

Which of the following are statistical questions?

\begin{enumerate}
\def\labelenumi{\arabic{enumi}.}
\tightlist
\item
  How many pages are in Harry Potter and the Philosopher's Stone?
\item
  How many hours per week do Year 11 students spend on homework?
\item
  What is the population of New Zealand?
\item
  What are the most popular music genres among teenagers in Auckland?
\item
  What colour is the school uniform?
\end{enumerate}

\begin{tcolorbox}[enhanced jigsaw, bottomtitle=1mm, opacityback=0, titlerule=0mm, leftrule=.75mm, opacitybacktitle=0.6, coltitle=black, toprule=.15mm, colbacktitle=quarto-callout-note-color!10!white, title=\textcolor{quarto-callout-note-color}{\faInfo}\hspace{0.5em}{Answers}, left=2mm, rightrule=.15mm, toptitle=1mm, colframe=quarto-callout-note-color-frame, breakable, arc=.35mm, bottomrule=.15mm, colback=white]

\begin{enumerate}
\def\labelenumi{\arabic{enumi}.}
\tightlist
\item
  ❌ Not statistical --- there is one fixed answer.
\item
  ✅ Statistical --- different students spend different amounts of time.
\item
  ❌ Not statistical --- there is one (changing but definite) answer.
\item
  ✅ Statistical --- different teenagers have different preferences.
\item
  ❌ Not statistical --- the uniform colour is fixed.
\end{enumerate}

\end{tcolorbox}

\section{Why context matters}\label{why-context-matters}

In AS91946, \textbf{context} is central. You will always work with data
that relates to real situations. When interpreting statistics, you must:

\begin{itemize}
\tightlist
\item
  \textbf{Name the variable} being measured
\item
  \textbf{State the units} (e.g.~centimetres, dollars, minutes)
\item
  \textbf{Refer to the population} the data came from
\item
  \textbf{Relate findings back} to the real-world situation
\end{itemize}

\begin{tcolorbox}[enhanced jigsaw, bottomtitle=1mm, opacityback=0, titlerule=0mm, leftrule=.75mm, opacitybacktitle=0.6, coltitle=black, toprule=.15mm, colbacktitle=quarto-callout-warning-color!10!white, title=\textcolor{quarto-callout-warning-color}{\faExclamationTriangle}\hspace{0.5em}{Common mistake}, left=2mm, rightrule=.15mm, toptitle=1mm, colframe=quarto-callout-warning-color-frame, breakable, arc=.35mm, bottomrule=.15mm, colback=white]

Saying ``the mean is 24'' is not enough. Say\\
\emph{``The mean age of survey respondents was 24 years.''}

\end{tcolorbox}

\section{Glossary}\label{glossary}

\begin{longtable}[]{@{}ll@{}}
\toprule\noalign{}
Term & Meaning \\
\midrule\noalign{}
\endhead
\bottomrule\noalign{}
\endlastfoot
\textbf{Data} & Values collected by observation or measurement \\
\textbf{Variable} & A characteristic that can take different values \\
\textbf{Population} & The entire group of interest \\
\textbf{Sample} & A subset of the population selected for study \\
\textbf{Statistic} & A numerical summary calculated from a sample \\
\textbf{Parameter} & A numerical summary of a whole population \\
\end{longtable}

\section{Check your understanding}\label{check-your-understanding}

\begin{enumerate}
\def\labelenumi{\arabic{enumi}.}
\tightlist
\item
  In your own words, explain what the PPDAC cycle is.
\item
  Give two examples of statistical questions you could investigate at
  your school.
\item
  A student says ``the average score is 65''. What is missing from this
  statement?
\item
  What is the difference between a population and a sample?
\end{enumerate}

\bookmarksetup{startatroot}

\chapter{Types of Data and Variables}\label{types-of-data-and-variables}

\bookmarksetup{startatroot}

\chapter{Types of Data and Variables}\label{sec-data-types}

Understanding the \textbf{type of data} you have is the first step in
choosing the right graph or summary statistic. The wrong choice can make
data hard to understand --- or even misleading.

\section{Categorical vs numerical
data}\label{categorical-vs-numerical-data}

\begin{figure}[H]

{\centering \pandocbounded{\includegraphics[keepaspectratio]{chapters/02-data-types_files/figure-latex/data-types-tree-1.pdf}}

}

\caption{Classification of data types.}

\end{figure}%

\section{Categorical data}\label{categorical-data}

\textbf{Categorical (qualitative) data} describes qualities or
categories. The values are labels, not numbers.

\subsection{Nominal data}\label{nominal-data}

Nominal data has \textbf{no natural order}. Categories are just names.

\textbf{Examples:} - Eye colour (brown, blue, green, grey) - Favourite
sport (rugby, netball, football, basketball) - Type of transport to
school (bus, car, walk, bike)

\subsection{Ordinal data}\label{ordinal-data}

Ordinal data has a \textbf{natural order}, but the gaps between
categories are not equal.

\textbf{Examples:} - Survey ratings (strongly disagree → disagree →
neutral → agree → strongly agree) - Year group (Year 9, Year 10, Year
11, Year 12, Year 13) - Clothing size (XS, S, M, L, XL)

\section{Numerical data}\label{numerical-data}

\textbf{Numerical (quantitative) data} represents quantities that can be
measured or counted.

\subsection{Discrete data}\label{discrete-data}

Discrete data can only take \textbf{specific, separate values} (usually
whole numbers from counting).

\textbf{Examples:} - Number of siblings - Number of text messages sent
per day - Goals scored in a rugby match

\subsection{Continuous data}\label{continuous-data}

Continuous data can take \textbf{any value within a range} (usually from
measuring).

\textbf{Examples:} - Height (cm) - Weight (kg) - Time to complete a race
(seconds) - Temperature (°C)

\section{Real data example}\label{real-data-example}

\begin{Shaded}
\begin{Highlighting}[]
\CommentTok{\# Simulated class survey data}
\FunctionTok{set.seed}\NormalTok{(}\DecValTok{7}\NormalTok{)}
\NormalTok{n }\OtherTok{\textless{}{-}} \DecValTok{10}
\NormalTok{survey }\OtherTok{\textless{}{-}} \FunctionTok{data.frame}\NormalTok{(}
  \AttributeTok{Student   =} \FunctionTok{paste0}\NormalTok{(}\StringTok{"S"}\NormalTok{, }\DecValTok{1}\SpecialCharTok{:}\NormalTok{n),}
  \AttributeTok{Gender    =} \FunctionTok{sample}\NormalTok{(}\FunctionTok{c}\NormalTok{(}\StringTok{"Male"}\NormalTok{, }\StringTok{"Female"}\NormalTok{, }\StringTok{"Non{-}binary"}\NormalTok{), n, }\AttributeTok{replace =} \ConstantTok{TRUE}\NormalTok{),}
  \AttributeTok{Height\_cm =} \FunctionTok{round}\NormalTok{(}\FunctionTok{rnorm}\NormalTok{(n, }\DecValTok{168}\NormalTok{, }\DecValTok{10}\NormalTok{), }\DecValTok{1}\NormalTok{),}
  \AttributeTok{Siblings  =} \FunctionTok{sample}\NormalTok{(}\DecValTok{0}\SpecialCharTok{:}\DecValTok{4}\NormalTok{, n, }\AttributeTok{replace =} \ConstantTok{TRUE}\NormalTok{),}
  \AttributeTok{Transport =} \FunctionTok{sample}\NormalTok{(}\FunctionTok{c}\NormalTok{(}\StringTok{"Bus"}\NormalTok{, }\StringTok{"Car"}\NormalTok{, }\StringTok{"Walk"}\NormalTok{, }\StringTok{"Bike"}\NormalTok{), n, }\AttributeTok{replace =} \ConstantTok{TRUE}\NormalTok{),}
  \AttributeTok{Rating    =} \FunctionTok{sample}\NormalTok{(}\FunctionTok{c}\NormalTok{(}\StringTok{"Poor"}\NormalTok{, }\StringTok{"Fair"}\NormalTok{, }\StringTok{"Good"}\NormalTok{, }\StringTok{"Excellent"}\NormalTok{), n, }\AttributeTok{replace =} \ConstantTok{TRUE}\NormalTok{)}
\NormalTok{)}

\NormalTok{knitr}\SpecialCharTok{::}\FunctionTok{kable}\NormalTok{(survey, }\AttributeTok{align =} \StringTok{"c"}\NormalTok{)}
\end{Highlighting}
\end{Shaded}

\begin{longtable}[]{@{}cccccc@{}}
\caption{Example student survey data --- identify the type of each
variable.}\tabularnewline
\toprule\noalign{}
Student & Gender & Height\_cm & Siblings & Transport & Rating \\
\midrule\noalign{}
\endfirsthead
\toprule\noalign{}
Student & Gender & Height\_cm & Siblings & Transport & Rating \\
\midrule\noalign{}
\endhead
\bottomrule\noalign{}
\endlastfoot
S1 & Female & 154.3 & 1 & Bus & Fair \\
S2 & Non-binary & 144.2 & 3 & Bike & Excellent \\
S3 & Non-binary & 163.2 & 2 & Bus & Excellent \\
S4 & Non-binary & 162.6 & 3 & Bus & Poor \\
S5 & Female & 181.2 & 0 & Car & Excellent \\
S6 & Non-binary & 152.9 & 2 & Bike & Fair \\
S7 & Female & 167.8 & 3 & Bus & Good \\
S8 & Female & 164.5 & 0 & Walk & Excellent \\
S9 & Non-binary & 161.7 & 2 & Bus & Fair \\
S10 & Female & 159.0 & 4 & Walk & Fair \\
\end{longtable}

\begin{tcolorbox}[enhanced jigsaw, bottomtitle=1mm, opacityback=0, titlerule=0mm, leftrule=.75mm, opacitybacktitle=0.6, coltitle=black, toprule=.15mm, colbacktitle=quarto-callout-tip-color!10!white, title=\textcolor{quarto-callout-tip-color}{\faLightbulb}\hspace{0.5em}{Identify the variable types}, left=2mm, rightrule=.15mm, toptitle=1mm, colframe=quarto-callout-tip-color-frame, breakable, arc=.35mm, bottomrule=.15mm, colback=white]

\begin{itemize}
\tightlist
\item
  \textbf{Gender} → Categorical, Nominal
\item
  \textbf{Height\_cm} → Numerical, Continuous
\item
  \textbf{Siblings} → Numerical, Discrete
\item
  \textbf{Transport} → Categorical, Nominal
\item
  \textbf{Rating} → Categorical, Ordinal
\end{itemize}

\end{tcolorbox}

\section{Choosing the right display}\label{choosing-the-right-display}

The type of data determines the best graph to use:

\begin{longtable}[]{@{}
  >{\raggedright\arraybackslash}p{(\linewidth - 2\tabcolsep) * \real{0.3793}}
  >{\raggedright\arraybackslash}p{(\linewidth - 2\tabcolsep) * \real{0.6207}}@{}}
\toprule\noalign{}
\begin{minipage}[b]{\linewidth}\raggedright
Data type
\end{minipage} & \begin{minipage}[b]{\linewidth}\raggedright
Suitable displays
\end{minipage} \\
\midrule\noalign{}
\endhead
\bottomrule\noalign{}
\endlastfoot
Categorical (nominal/ordinal) & Bar chart, pie chart, dot plot \\
Numerical (discrete) & Dot plot, bar chart, stem-and-leaf \\
Numerical (continuous) & Histogram, box plot, dot plot, stem-and-leaf \\
Two numerical variables & Scatter plot \\
\end{longtable}

\section{Glossary}\label{glossary-1}

\begin{longtable}[]{@{}ll@{}}
\toprule\noalign{}
Term & Meaning \\
\midrule\noalign{}
\endhead
\bottomrule\noalign{}
\endlastfoot
\textbf{Variable} & A characteristic that varies between individuals \\
\textbf{Categorical} & Data that falls into named groups or
categories \\
\textbf{Nominal} & Categorical with no natural order \\
\textbf{Ordinal} & Categorical with a natural order \\
\textbf{Numerical} & Data that can be measured or counted \\
\textbf{Discrete} & Numerical data that can only take specific values \\
\textbf{Continuous} & Numerical data that can take any value in a
range \\
\end{longtable}

\section{Check your understanding}\label{check-your-understanding-1}

Classify each variable as \textbf{categorical (nominal/ordinal)} or
\textbf{numerical (discrete/continuous)}:

\begin{enumerate}
\def\labelenumi{\arabic{enumi}.}
\tightlist
\item
  Shoe size
\item
  Favourite subject at school
\item
  Time spent sleeping (hours per night)
\item
  Number of books read last month
\item
  Blood type
\item
  Temperature in a classroom
\item
  Year level
\item
  Distance from home to school (km)
\end{enumerate}

\begin{tcolorbox}[enhanced jigsaw, bottomtitle=1mm, opacityback=0, titlerule=0mm, leftrule=.75mm, opacitybacktitle=0.6, coltitle=black, toprule=.15mm, colbacktitle=quarto-callout-note-color!10!white, title=\textcolor{quarto-callout-note-color}{\faInfo}\hspace{0.5em}{Answers}, left=2mm, rightrule=.15mm, toptitle=1mm, colframe=quarto-callout-note-color-frame, breakable, arc=.35mm, bottomrule=.15mm, colback=white]

\begin{enumerate}
\def\labelenumi{\arabic{enumi}.}
\tightlist
\item
  Numerical, discrete (or ordinal --- shoe sizes have a natural order
  but are not continuous)
\item
  Categorical, nominal
\item
  Numerical, continuous
\item
  Numerical, discrete
\item
  Categorical, nominal
\item
  Numerical, continuous
\item
  Categorical, ordinal (or numerical discrete depending on context)
\item
  Numerical, continuous
\end{enumerate}

\end{tcolorbox}

\bookmarksetup{startatroot}

\chapter{Statistical Graphs and
Displays}\label{statistical-graphs-and-displays}

\bookmarksetup{startatroot}

\chapter{Statistical Graphs and Displays}\label{sec-graphs}

A well-chosen graph makes patterns in data immediately visible. In this
chapter you will learn to \textbf{create}, \textbf{read}, and
\textbf{critically interpret} the most common statistical displays.

\section{Setting up our data}\label{setting-up-our-data}

Throughout this chapter we use a simulated survey of 60 Year 11 students
across a New Zealand school. The data includes variables on transport,
height, study time, and maths scores.

\begin{Shaded}
\begin{Highlighting}[]
\FunctionTok{set.seed}\NormalTok{(}\DecValTok{2024}\NormalTok{)}
\NormalTok{n }\OtherTok{\textless{}{-}} \DecValTok{60}

\NormalTok{students }\OtherTok{\textless{}{-}} \FunctionTok{data.frame}\NormalTok{(}
  \AttributeTok{transport  =} \FunctionTok{sample}\NormalTok{(}\FunctionTok{c}\NormalTok{(}\StringTok{"Bus"}\NormalTok{, }\StringTok{"Car"}\NormalTok{, }\StringTok{"Walk"}\NormalTok{, }\StringTok{"Bike"}\NormalTok{), n,}
                      \AttributeTok{replace =} \ConstantTok{TRUE}\NormalTok{, }\AttributeTok{prob =} \FunctionTok{c}\NormalTok{(}\FloatTok{0.4}\NormalTok{, }\FloatTok{0.25}\NormalTok{, }\FloatTok{0.25}\NormalTok{, }\FloatTok{0.1}\NormalTok{)),}
  \AttributeTok{height\_cm  =} \FunctionTok{round}\NormalTok{(}\FunctionTok{rnorm}\NormalTok{(n, }\AttributeTok{mean =} \DecValTok{167}\NormalTok{, }\AttributeTok{sd =} \DecValTok{9}\NormalTok{), }\DecValTok{1}\NormalTok{),}
  \AttributeTok{study\_hrs  =} \FunctionTok{round}\NormalTok{(}\FunctionTok{pmax}\NormalTok{(}\DecValTok{0}\NormalTok{, }\FunctionTok{rnorm}\NormalTok{(n, }\AttributeTok{mean =} \FloatTok{2.5}\NormalTok{, }\AttributeTok{sd =} \FloatTok{1.2}\NormalTok{)), }\DecValTok{1}\NormalTok{),}
  \AttributeTok{maths\_score =} \FunctionTok{round}\NormalTok{(}\FunctionTok{pmin}\NormalTok{(}\DecValTok{100}\NormalTok{, }\FunctionTok{pmax}\NormalTok{(}\DecValTok{0}\NormalTok{, }\FunctionTok{rnorm}\NormalTok{(n, }\AttributeTok{mean =} \DecValTok{65}\NormalTok{, }\AttributeTok{sd =} \DecValTok{15}\NormalTok{))))}
\NormalTok{)}
\end{Highlighting}
\end{Shaded}

\section{Bar charts}\label{bar-charts}

Bar charts display the \textbf{frequency} (count) or \textbf{relative
frequency} (proportion) of categorical data.

\begin{Shaded}
\begin{Highlighting}[]
\NormalTok{transport\_counts }\OtherTok{\textless{}{-}} \FunctionTok{table}\NormalTok{(students}\SpecialCharTok{$}\NormalTok{transport)}

\FunctionTok{barplot}\NormalTok{(}
  \FunctionTok{sort}\NormalTok{(transport\_counts, }\AttributeTok{decreasing =} \ConstantTok{TRUE}\NormalTok{),}
  \AttributeTok{col  =} \FunctionTok{c}\NormalTok{(}\StringTok{"\#4dac26"}\NormalTok{, }\StringTok{"\#b8e186"}\NormalTok{, }\StringTok{"\#f1b6da"}\NormalTok{, }\StringTok{"\#d7191c"}\NormalTok{),}
  \AttributeTok{main =} \StringTok{"How Year 11 Students Travel to School"}\NormalTok{,}
  \AttributeTok{xlab =} \StringTok{"Mode of Transport"}\NormalTok{,}
  \AttributeTok{ylab =} \StringTok{"Number of Students"}\NormalTok{,}
  \AttributeTok{border =} \StringTok{"white"}\NormalTok{,}
  \AttributeTok{las =} \DecValTok{1}
\NormalTok{)}
\end{Highlighting}
\end{Shaded}

\begin{figure}[H]

\centering{

\pandocbounded{\includegraphics[keepaspectratio]{chapters/03-graphs_files/figure-latex/fig-bar-transport-1.pdf}}

}

\caption{\label{fig-bar-transport}Bar chart showing how Year 11 students
travel to school.}

\end{figure}%

\subsection{Reading a bar chart}\label{reading-a-bar-chart}

\begin{tcolorbox}[enhanced jigsaw, bottomtitle=1mm, opacityback=0, titlerule=0mm, leftrule=.75mm, opacitybacktitle=0.6, coltitle=black, toprule=.15mm, colbacktitle=quarto-callout-tip-color!10!white, title=\textcolor{quarto-callout-tip-color}{\faLightbulb}\hspace{0.5em}{Tip}, left=2mm, rightrule=.15mm, toptitle=1mm, colframe=quarto-callout-tip-color-frame, breakable, arc=.35mm, bottomrule=.15mm, colback=white]

When describing a bar chart, always comment on:

\begin{enumerate}
\def\labelenumi{\arabic{enumi}.}
\tightlist
\item
  \textbf{The most and least common} categories
\item
  \textbf{The overall pattern} (are categories roughly equal, or is one
  dominant?)
\item
  \textbf{The context} --- what does this tell us about the group?
\end{enumerate}

\end{tcolorbox}

\textbf{Sample answer:} ``Bus is the most common way students travel to
school, with over 20 students using this method. Bike is the least
common. More students use public or active transport (bus, walk, bike)
than passive transport (car).''

\section{Dot plots}\label{dot-plots}

Dot plots are useful for \textbf{small numerical datasets}. Each data
point is shown as a dot above a number line.

\begin{Shaded}
\begin{Highlighting}[]
\CommentTok{\# Rounded to nearest 0.5 for clearer dot plot}
\NormalTok{study\_rounded }\OtherTok{\textless{}{-}} \FunctionTok{round}\NormalTok{(students}\SpecialCharTok{$}\NormalTok{study\_hrs }\SpecialCharTok{*} \DecValTok{2}\NormalTok{) }\SpecialCharTok{/} \DecValTok{2}

\FunctionTok{stripchart}\NormalTok{(}
\NormalTok{  study\_rounded,}
  \AttributeTok{method =} \StringTok{"stack"}\NormalTok{,}
  \AttributeTok{pch    =} \DecValTok{19}\NormalTok{,}
  \AttributeTok{col    =} \StringTok{"\#0571b0"}\NormalTok{,}
  \AttributeTok{cex    =} \FloatTok{1.4}\NormalTok{,}
  \AttributeTok{xlab   =} \StringTok{"Daily Study Hours"}\NormalTok{,}
  \AttributeTok{main   =} \StringTok{"Daily Study Hours — Year 11 Students"}\NormalTok{,}
  \AttributeTok{offset =} \FloatTok{0.4}
\NormalTok{)}
\end{Highlighting}
\end{Shaded}

\begin{figure}[H]

\centering{

\pandocbounded{\includegraphics[keepaspectratio]{chapters/03-graphs_files/figure-latex/fig-dotplot-study-1.pdf}}

}

\caption{\label{fig-dotplot-study}Dot plot of daily study hours for Year
11 students.}

\end{figure}%

\section{Histograms}\label{histograms}

A histogram shows the \textbf{distribution of continuous numerical data}
by grouping values into intervals (bins).

\begin{Shaded}
\begin{Highlighting}[]
\FunctionTok{hist}\NormalTok{(}
\NormalTok{  students}\SpecialCharTok{$}\NormalTok{height\_cm,}
  \AttributeTok{breaks =} \DecValTok{10}\NormalTok{,}
  \AttributeTok{col    =} \StringTok{"\#2c7bb6"}\NormalTok{,}
  \AttributeTok{border =} \StringTok{"white"}\NormalTok{,}
  \AttributeTok{main   =} \StringTok{"Distribution of Student Heights"}\NormalTok{,}
  \AttributeTok{xlab   =} \StringTok{"Height (cm)"}\NormalTok{,}
  \AttributeTok{ylab   =} \StringTok{"Frequency"}\NormalTok{,}
  \AttributeTok{las    =} \DecValTok{1}
\NormalTok{)}
\end{Highlighting}
\end{Shaded}

\begin{figure}[H]

\centering{

\pandocbounded{\includegraphics[keepaspectratio]{chapters/03-graphs_files/figure-latex/fig-histogram-height-1.pdf}}

}

\caption{\label{fig-histogram-height}Histogram of student heights (cm).}

\end{figure}%

\subsection{Describing a histogram ---
SOCS}\label{describing-a-histogram-socs}

Use \textbf{SOCS} to describe a distribution:

\begin{longtable}[]{@{}
  >{\raggedright\arraybackslash}p{(\linewidth - 2\tabcolsep) * \real{0.3200}}
  >{\raggedright\arraybackslash}p{(\linewidth - 2\tabcolsep) * \real{0.6800}}@{}}
\toprule\noalign{}
\begin{minipage}[b]{\linewidth}\raggedright
Letter
\end{minipage} & \begin{minipage}[b]{\linewidth}\raggedright
What to describe
\end{minipage} \\
\midrule\noalign{}
\endhead
\bottomrule\noalign{}
\endlastfoot
\textbf{S}hape & Symmetric, skewed left, skewed right, uniform,
bimodal \\
\textbf{O}utliers & Any unusual values far from the main group \\
\textbf{C}entre & Where the middle of the data is (mean or median) \\
\textbf{S}pread & How variable the data is (range, IQR, standard
deviation) \\
\end{longtable}

\begin{Shaded}
\begin{Highlighting}[]
\FunctionTok{par}\NormalTok{(}\AttributeTok{mfrow =} \FunctionTok{c}\NormalTok{(}\DecValTok{1}\NormalTok{, }\DecValTok{3}\NormalTok{))}

\CommentTok{\# Symmetric}
\FunctionTok{set.seed}\NormalTok{(}\DecValTok{1}\NormalTok{)}
\FunctionTok{hist}\NormalTok{(}\FunctionTok{rnorm}\NormalTok{(}\DecValTok{500}\NormalTok{, }\DecValTok{50}\NormalTok{, }\DecValTok{10}\NormalTok{), }\AttributeTok{breaks =} \DecValTok{15}\NormalTok{, }\AttributeTok{col =} \StringTok{"\#4dac26"}\NormalTok{, }\AttributeTok{border =} \StringTok{"white"}\NormalTok{,}
     \AttributeTok{main =} \StringTok{"Symmetric"}\NormalTok{, }\AttributeTok{xlab =} \StringTok{""}\NormalTok{, }\AttributeTok{ylab =} \StringTok{"Frequency"}\NormalTok{, }\AttributeTok{las =} \DecValTok{1}\NormalTok{)}

\CommentTok{\# Right{-}skewed}
\FunctionTok{hist}\NormalTok{(}\FunctionTok{rexp}\NormalTok{(}\DecValTok{500}\NormalTok{, }\FloatTok{0.1}\NormalTok{), }\AttributeTok{breaks =} \DecValTok{15}\NormalTok{, }\AttributeTok{col =} \StringTok{"\#f4a582"}\NormalTok{, }\AttributeTok{border =} \StringTok{"white"}\NormalTok{,}
     \AttributeTok{main =} \StringTok{"Right (positive) skew"}\NormalTok{, }\AttributeTok{xlab =} \StringTok{""}\NormalTok{, }\AttributeTok{ylab =} \StringTok{""}\NormalTok{, }\AttributeTok{las =} \DecValTok{1}\NormalTok{)}

\CommentTok{\# Left{-}skewed}
\FunctionTok{hist}\NormalTok{(}\DecValTok{100} \SpecialCharTok{{-}} \FunctionTok{rexp}\NormalTok{(}\DecValTok{500}\NormalTok{, }\FloatTok{0.1}\NormalTok{), }\AttributeTok{breaks =} \DecValTok{15}\NormalTok{, }\AttributeTok{col =} \StringTok{"\#92c5de"}\NormalTok{, }\AttributeTok{border =} \StringTok{"white"}\NormalTok{,}
     \AttributeTok{main =} \StringTok{"Left (negative) skew"}\NormalTok{, }\AttributeTok{xlab =} \StringTok{""}\NormalTok{, }\AttributeTok{ylab =} \StringTok{""}\NormalTok{, }\AttributeTok{las =} \DecValTok{1}\NormalTok{)}

\FunctionTok{par}\NormalTok{(}\AttributeTok{mfrow =} \FunctionTok{c}\NormalTok{(}\DecValTok{1}\NormalTok{, }\DecValTok{1}\NormalTok{))}
\end{Highlighting}
\end{Shaded}

\begin{figure}[H]

\centering{

\pandocbounded{\includegraphics[keepaspectratio]{chapters/03-graphs_files/figure-latex/fig-shape-examples-1.pdf}}

}

\caption{\label{fig-shape-examples}Examples of different distribution
shapes.}

\end{figure}%

\section{Box plots (box-and-whisker
plots)}\label{box-plots-box-and-whisker-plots}

A box plot summarises a \textbf{five-number summary}: Minimum, Q1,
Median, Q3, Maximum.

\begin{Shaded}
\begin{Highlighting}[]
\FunctionTok{boxplot}\NormalTok{(}
\NormalTok{  maths\_score }\SpecialCharTok{\textasciitilde{}}\NormalTok{ transport,}
  \AttributeTok{data   =}\NormalTok{ students,}
  \AttributeTok{col    =} \FunctionTok{c}\NormalTok{(}\StringTok{"\#4dac26"}\NormalTok{, }\StringTok{"\#f4a582"}\NormalTok{, }\StringTok{"\#92c5de"}\NormalTok{, }\StringTok{"\#d7191c"}\NormalTok{),}
  \AttributeTok{border =} \StringTok{"\#333333"}\NormalTok{,}
  \AttributeTok{main   =} \StringTok{"Maths Scores by Transport Mode"}\NormalTok{,}
  \AttributeTok{xlab   =} \StringTok{"Mode of Transport"}\NormalTok{,}
  \AttributeTok{ylab   =} \StringTok{"Maths Score"}\NormalTok{,}
  \AttributeTok{las    =} \DecValTok{1}
\NormalTok{)}
\end{Highlighting}
\end{Shaded}

\begin{figure}[H]

\centering{

\pandocbounded{\includegraphics[keepaspectratio]{chapters/03-graphs_files/figure-latex/fig-boxplot-1.pdf}}

}

\caption{\label{fig-boxplot}Box plot of maths scores by transport type.}

\end{figure}%

\subsection{Parts of a box plot}\label{parts-of-a-box-plot}

\begin{figure}

\centering{

\pandocbounded{\includegraphics[keepaspectratio]{chapters/03-graphs_files/figure-latex/fig-boxplot-anatomy-1.pdf}}

}

\caption{\label{fig-boxplot-anatomy}Anatomy of a box plot.}

\end{figure}%

\section{Stem-and-leaf plots}\label{stem-and-leaf-plots}

Stem-and-leaf plots show the actual data values and the shape of the
distribution simultaneously. They work well for \textbf{small datasets}
(up to \textasciitilde50 values).

\begin{Shaded}
\begin{Highlighting}[]
\CommentTok{\# Show a stem{-}and{-}leaf for a subset of maths scores}
\FunctionTok{cat}\NormalTok{(}\StringTok{"Stem{-}and{-}leaf plot of maths scores (first 30 students):}\SpecialCharTok{\textbackslash{}n\textbackslash{}n}\StringTok{"}\NormalTok{)}
\CommentTok{\#\textgreater{} Stem{-}and{-}leaf plot of maths scores (first 30 students):}
\FunctionTok{stem}\NormalTok{(students}\SpecialCharTok{$}\NormalTok{maths\_score[}\DecValTok{1}\SpecialCharTok{:}\DecValTok{30}\NormalTok{])}
\CommentTok{\#\textgreater{} }
\CommentTok{\#\textgreater{}   The decimal point is 1 digit(s) to the right of the |}
\CommentTok{\#\textgreater{} }
\CommentTok{\#\textgreater{}    2 | 1}
\CommentTok{\#\textgreater{}    3 | 9}
\CommentTok{\#\textgreater{}    4 | 8}
\CommentTok{\#\textgreater{}    5 | 3468999}
\CommentTok{\#\textgreater{}    6 | 04557789}
\CommentTok{\#\textgreater{}    7 | 0169}
\CommentTok{\#\textgreater{}    8 | 23448}
\CommentTok{\#\textgreater{}    9 | 18}
\CommentTok{\#\textgreater{}   10 | 0}
\end{Highlighting}
\end{Shaded}

\section{Scatter plots}\label{scatter-plots}

Scatter plots show the \textbf{relationship between two numerical
variables}.

\begin{Shaded}
\begin{Highlighting}[]
\FunctionTok{plot}\NormalTok{(}
\NormalTok{  students}\SpecialCharTok{$}\NormalTok{study\_hrs, students}\SpecialCharTok{$}\NormalTok{maths\_score,}
  \AttributeTok{pch  =} \DecValTok{19}\NormalTok{,}
  \AttributeTok{col  =} \FunctionTok{adjustcolor}\NormalTok{(}\StringTok{"\#0571b0"}\NormalTok{, }\AttributeTok{alpha.f =} \FloatTok{0.7}\NormalTok{),}
  \AttributeTok{xlab =} \StringTok{"Daily Study Hours"}\NormalTok{,}
  \AttributeTok{ylab =} \StringTok{"Maths Score"}\NormalTok{,}
  \AttributeTok{main =} \StringTok{"Study Hours vs Maths Score — Year 11 Students"}
\NormalTok{)}
\FunctionTok{abline}\NormalTok{(}\FunctionTok{lm}\NormalTok{(maths\_score }\SpecialCharTok{\textasciitilde{}}\NormalTok{ study\_hrs, }\AttributeTok{data =}\NormalTok{ students),}
       \AttributeTok{col =} \StringTok{"\#d7191c"}\NormalTok{, }\AttributeTok{lwd =} \DecValTok{2}\NormalTok{)}
\end{Highlighting}
\end{Shaded}

\begin{figure}[H]

\centering{

\pandocbounded{\includegraphics[keepaspectratio]{chapters/03-graphs_files/figure-latex/fig-scatter-1.pdf}}

}

\caption{\label{fig-scatter}Scatter plot of study hours vs maths score.}

\end{figure}%

\subsection{Describing a scatter plot}\label{describing-a-scatter-plot}

Comment on:

\begin{enumerate}
\def\labelenumi{\arabic{enumi}.}
\tightlist
\item
  \textbf{Direction} --- positive, negative, or no association
\item
  \textbf{Form} --- linear or non-linear
\item
  \textbf{Strength} --- strong, moderate, or weak
\item
  \textbf{Outliers} --- any unusual points
\end{enumerate}

\begin{tcolorbox}[enhanced jigsaw, bottomtitle=1mm, opacityback=0, titlerule=0mm, leftrule=.75mm, opacitybacktitle=0.6, coltitle=black, toprule=.15mm, colbacktitle=quarto-callout-warning-color!10!white, title=\textcolor{quarto-callout-warning-color}{\faExclamationTriangle}\hspace{0.5em}{Correlation ≠ Causation}, left=2mm, rightrule=.15mm, toptitle=1mm, colframe=quarto-callout-warning-color-frame, breakable, arc=.35mm, bottomrule=.15mm, colback=white]

Just because two variables are related \textbf{does not} mean one causes
the other. Always think about whether an association makes sense in
context.

\end{tcolorbox}

\section{Pie charts}\label{pie-charts}

Pie charts show \textbf{proportions} of categorical data. Use them
carefully --- bar charts are usually easier to read precisely.

\begin{Shaded}
\begin{Highlighting}[]
\FunctionTok{pie}\NormalTok{(}
\NormalTok{  transport\_counts,}
  \AttributeTok{col    =} \FunctionTok{c}\NormalTok{(}\StringTok{"\#4dac26"}\NormalTok{, }\StringTok{"\#f4a582"}\NormalTok{, }\StringTok{"\#92c5de"}\NormalTok{, }\StringTok{"\#d7191c"}\NormalTok{),}
  \AttributeTok{main   =} \StringTok{"Transport Modes — Year 11 Students"}\NormalTok{,}
  \AttributeTok{labels =} \FunctionTok{paste0}\NormalTok{(}\FunctionTok{names}\NormalTok{(transport\_counts), }\StringTok{"}\SpecialCharTok{\textbackslash{}n}\StringTok{("}\NormalTok{,}
                  \FunctionTok{round}\NormalTok{(}\DecValTok{100} \SpecialCharTok{*}\NormalTok{ transport\_counts }\SpecialCharTok{/} \FunctionTok{sum}\NormalTok{(transport\_counts)), }\StringTok{"\%)"}\NormalTok{)}
\NormalTok{)}
\end{Highlighting}
\end{Shaded}

\begin{figure}[H]

\centering{

\pandocbounded{\includegraphics[keepaspectratio]{chapters/03-graphs_files/figure-latex/fig-pie-1.pdf}}

}

\caption{\label{fig-pie}Pie chart of transport modes.}

\end{figure}%

\section{Choosing the right graph}\label{choosing-the-right-graph}

\begin{longtable}[]{@{}
  >{\raggedright\arraybackslash}p{(\linewidth - 2\tabcolsep) * \real{0.5000}}
  >{\raggedright\arraybackslash}p{(\linewidth - 2\tabcolsep) * \real{0.5000}}@{}}
\toprule\noalign{}
\begin{minipage}[b]{\linewidth}\raggedright
Situation
\end{minipage} & \begin{minipage}[b]{\linewidth}\raggedright
Best graph
\end{minipage} \\
\midrule\noalign{}
\endhead
\bottomrule\noalign{}
\endlastfoot
One categorical variable & Bar chart, pie chart \\
One numerical variable (small n) & Dot plot, stem-and-leaf \\
One numerical variable (large n) & Histogram, box plot \\
Comparing groups numerically & Side-by-side box plots, back-to-back
stem-and-leaf \\
Two numerical variables & Scatter plot \\
\end{longtable}

\section{Check your understanding}\label{check-your-understanding-2}

\begin{enumerate}
\def\labelenumi{\arabic{enumi}.}
\item
  A student recorded the heights of 25 classmates. Which graph type is
  most appropriate to display this data? Explain your choice.
\item
  Describe the scatter plot in Figure~\ref{fig-scatter} using the four
  features (direction, form, strength, outliers).
\item
  A bar chart shows the following transport counts: Bus = 24, Car = 15,
  Walk = 15, Bike = 6. What percentage of students walk to school?
\item
  True or false: ``A right-skewed distribution has a longer tail on the
  left.'' Explain.
\end{enumerate}

\begin{tcolorbox}[enhanced jigsaw, bottomtitle=1mm, opacityback=0, titlerule=0mm, leftrule=.75mm, opacitybacktitle=0.6, coltitle=black, toprule=.15mm, colbacktitle=quarto-callout-note-color!10!white, title=\textcolor{quarto-callout-note-color}{\faInfo}\hspace{0.5em}{Answers}, left=2mm, rightrule=.15mm, toptitle=1mm, colframe=quarto-callout-note-color-frame, breakable, arc=.35mm, bottomrule=.15mm, colback=white]

\begin{enumerate}
\def\labelenumi{\arabic{enumi}.}
\tightlist
\item
  A histogram or dot plot --- height is numerical and continuous.
\item
  Positive direction (more study → higher score), roughly linear form,
  moderate strength, possibly one or two outliers.
\item
  Walk = 15/60 = 25\%.
\item
  False --- a right-skewed (positively skewed) distribution has a longer
  tail on the RIGHT.
\end{enumerate}

\end{tcolorbox}

\bookmarksetup{startatroot}

\chapter{Measures of Centre and
Spread}\label{measures-of-centre-and-spread}

\bookmarksetup{startatroot}

\chapter{Measures of Centre and Spread}\label{sec-measures}

Once you have a graph, you need \textbf{numerical summaries} to describe
and compare distributions precisely. The two key things to describe are:

\begin{itemize}
\tightlist
\item
  \textbf{Centre} --- where is the ``middle'' of the data?
\item
  \textbf{Spread} --- how variable is the data?
\end{itemize}

\section{Measures of centre}\label{measures-of-centre}

\subsection{The mean}\label{the-mean}

The \textbf{mean} (average) is calculated by adding all values and
dividing by the number of values:

\[\bar{x} = \frac{\sum x_i}{n}\]

\begin{Shaded}
\begin{Highlighting}[]
\FunctionTok{set.seed}\NormalTok{(}\DecValTok{2024}\NormalTok{)}
\NormalTok{scores }\OtherTok{\textless{}{-}} \FunctionTok{c}\NormalTok{(}\DecValTok{55}\NormalTok{, }\DecValTok{62}\NormalTok{, }\DecValTok{68}\NormalTok{, }\DecValTok{70}\NormalTok{, }\DecValTok{72}\NormalTok{, }\DecValTok{75}\NormalTok{, }\DecValTok{78}\NormalTok{, }\DecValTok{80}\NormalTok{, }\DecValTok{83}\NormalTok{, }\DecValTok{85}\NormalTok{,}
            \DecValTok{88}\NormalTok{, }\DecValTok{60}\NormalTok{, }\DecValTok{91}\NormalTok{, }\DecValTok{74}\NormalTok{, }\DecValTok{66}\NormalTok{, }\DecValTok{79}\NormalTok{, }\DecValTok{52}\NormalTok{, }\DecValTok{95}\NormalTok{, }\DecValTok{73}\NormalTok{, }\DecValTok{71}\NormalTok{)}
\NormalTok{n }\OtherTok{\textless{}{-}} \FunctionTok{length}\NormalTok{(scores)}

\FunctionTok{cat}\NormalTok{(}\StringTok{"Scores:"}\NormalTok{, }\FunctionTok{sort}\NormalTok{(scores), }\StringTok{"}\SpecialCharTok{\textbackslash{}n}\StringTok{"}\NormalTok{)}
\CommentTok{\#\textgreater{} Scores: 52 55 60 62 66 68 70 71 72 73 74 75 78 79 80 83 85 88 91 95}
\FunctionTok{cat}\NormalTok{(}\StringTok{"Mean:  "}\NormalTok{, }\FunctionTok{round}\NormalTok{(}\FunctionTok{mean}\NormalTok{(scores), }\DecValTok{2}\NormalTok{), }\StringTok{"}\SpecialCharTok{\textbackslash{}n}\StringTok{"}\NormalTok{)}
\CommentTok{\#\textgreater{} Mean:   73.85}
\FunctionTok{cat}\NormalTok{(}\StringTok{"n =    "}\NormalTok{, n, }\StringTok{"}\SpecialCharTok{\textbackslash{}n}\StringTok{"}\NormalTok{)}
\CommentTok{\#\textgreater{} n =     20}
\end{Highlighting}
\end{Shaded}

\subsection{The median}\label{the-median}

The \textbf{median} is the middle value when data is sorted from
smallest to largest. It is not affected by outliers.

\begin{itemize}
\tightlist
\item
  If \emph{n} is \textbf{odd}: median = middle value
\item
  If \emph{n} is \textbf{even}: median = average of the two middle
  values
\end{itemize}

\begin{Shaded}
\begin{Highlighting}[]
\NormalTok{sorted\_scores }\OtherTok{\textless{}{-}} \FunctionTok{sort}\NormalTok{(scores)}
\FunctionTok{cat}\NormalTok{(}\StringTok{"Sorted scores:}\SpecialCharTok{\textbackslash{}n}\StringTok{"}\NormalTok{, sorted\_scores, }\StringTok{"}\SpecialCharTok{\textbackslash{}n\textbackslash{}n}\StringTok{"}\NormalTok{)}
\CommentTok{\#\textgreater{} Sorted scores:}
\CommentTok{\#\textgreater{}  52 55 60 62 66 68 70 71 72 73 74 75 78 79 80 83 85 88 91 95}
\FunctionTok{cat}\NormalTok{(}\StringTok{"n = 20 (even) → median = mean of 10th and 11th values}\SpecialCharTok{\textbackslash{}n}\StringTok{"}\NormalTok{)}
\CommentTok{\#\textgreater{} n = 20 (even) → median = mean of 10th and 11th values}
\FunctionTok{cat}\NormalTok{(}\StringTok{"10th value:"}\NormalTok{, sorted\_scores[}\DecValTok{10}\NormalTok{], }\StringTok{"}\SpecialCharTok{\textbackslash{}n}\StringTok{"}\NormalTok{)}
\CommentTok{\#\textgreater{} 10th value: 73}
\FunctionTok{cat}\NormalTok{(}\StringTok{"11th value:"}\NormalTok{, sorted\_scores[}\DecValTok{11}\NormalTok{], }\StringTok{"}\SpecialCharTok{\textbackslash{}n}\StringTok{"}\NormalTok{)}
\CommentTok{\#\textgreater{} 11th value: 74}
\FunctionTok{cat}\NormalTok{(}\StringTok{"Median ="}\NormalTok{, }\FunctionTok{median}\NormalTok{(scores), }\StringTok{"}\SpecialCharTok{\textbackslash{}n}\StringTok{"}\NormalTok{)}
\CommentTok{\#\textgreater{} Median = 73.5}
\end{Highlighting}
\end{Shaded}

\subsection{The mode}\label{the-mode}

The \textbf{mode} is the most frequently occurring value. Data can have:
- \textbf{No mode} (all values appear once) - \textbf{One mode}
(unimodal) - \textbf{Two modes} (bimodal) - \textbf{More than two modes}
(multimodal)

\begin{Shaded}
\begin{Highlighting}[]
\CommentTok{\# Custom mode function (base R has no built{-}in mode for numeric)}
\NormalTok{stat\_mode }\OtherTok{\textless{}{-}} \ControlFlowTok{function}\NormalTok{(x) \{}
\NormalTok{  tbl }\OtherTok{\textless{}{-}} \FunctionTok{table}\NormalTok{(x)}
  \FunctionTok{as.numeric}\NormalTok{(}\FunctionTok{names}\NormalTok{(tbl)[tbl }\SpecialCharTok{==} \FunctionTok{max}\NormalTok{(tbl)])}
\NormalTok{\}}

\NormalTok{test\_data }\OtherTok{\textless{}{-}} \FunctionTok{c}\NormalTok{(}\DecValTok{4}\NormalTok{, }\DecValTok{7}\NormalTok{, }\DecValTok{2}\NormalTok{, }\DecValTok{7}\NormalTok{, }\DecValTok{9}\NormalTok{, }\DecValTok{7}\NormalTok{, }\DecValTok{3}\NormalTok{, }\DecValTok{4}\NormalTok{, }\DecValTok{5}\NormalTok{, }\DecValTok{4}\NormalTok{)}
\FunctionTok{cat}\NormalTok{(}\StringTok{"Data:"}\NormalTok{, test\_data, }\StringTok{"}\SpecialCharTok{\textbackslash{}n}\StringTok{"}\NormalTok{)}
\CommentTok{\#\textgreater{} Data: 4 7 2 7 9 7 3 4 5 4}
\FunctionTok{cat}\NormalTok{(}\StringTok{"Mode(s):"}\NormalTok{, }\FunctionTok{stat\_mode}\NormalTok{(test\_data), }\StringTok{"}\SpecialCharTok{\textbackslash{}n}\StringTok{"}\NormalTok{)}
\CommentTok{\#\textgreater{} Mode(s): 4 7}
\end{Highlighting}
\end{Shaded}

\subsection{When to use each measure of
centre}\label{when-to-use-each-measure-of-centre}

\begin{longtable}[]{@{}lll@{}}
\toprule\noalign{}
Measure & Best used when\ldots{} & Affected by outliers? \\
\midrule\noalign{}
\endhead
\bottomrule\noalign{}
\endlastfoot
\textbf{Mean} & Data is roughly symmetric, no outliers & Yes \\
\textbf{Median} & Data is skewed or has outliers & No \\
\textbf{Mode} & Data is categorical; or finding most common value &
No \\
\end{longtable}

\begin{Shaded}
\begin{Highlighting}[]
\FunctionTok{par}\NormalTok{(}\AttributeTok{mfrow =} \FunctionTok{c}\NormalTok{(}\DecValTok{1}\NormalTok{, }\DecValTok{2}\NormalTok{))}

\CommentTok{\# Without outlier}
\NormalTok{x1 }\OtherTok{\textless{}{-}} \FunctionTok{c}\NormalTok{(}\DecValTok{10}\NormalTok{, }\DecValTok{12}\NormalTok{, }\DecValTok{13}\NormalTok{, }\DecValTok{14}\NormalTok{, }\DecValTok{15}\NormalTok{, }\DecValTok{16}\NormalTok{, }\DecValTok{17}\NormalTok{, }\DecValTok{18}\NormalTok{, }\DecValTok{19}\NormalTok{, }\DecValTok{20}\NormalTok{)}
\FunctionTok{hist}\NormalTok{(x1, }\AttributeTok{col =} \StringTok{"\#4dac26"}\NormalTok{, }\AttributeTok{border =} \StringTok{"white"}\NormalTok{, }\AttributeTok{main =} \StringTok{"No Outlier"}\NormalTok{,}
     \AttributeTok{xlab =} \StringTok{"Value"}\NormalTok{, }\AttributeTok{ylab =} \StringTok{"Frequency"}\NormalTok{, }\AttributeTok{xlim =} \FunctionTok{c}\NormalTok{(}\DecValTok{0}\NormalTok{, }\DecValTok{80}\NormalTok{), }\AttributeTok{breaks =} \DecValTok{5}\NormalTok{)}
\FunctionTok{abline}\NormalTok{(}\AttributeTok{v =} \FunctionTok{mean}\NormalTok{(x1),   }\AttributeTok{col =} \StringTok{"\#d7191c"}\NormalTok{, }\AttributeTok{lwd =} \DecValTok{2}\NormalTok{, }\AttributeTok{lty =} \DecValTok{1}\NormalTok{)}
\FunctionTok{abline}\NormalTok{(}\AttributeTok{v =} \FunctionTok{median}\NormalTok{(x1), }\AttributeTok{col =} \StringTok{"\#0571b0"}\NormalTok{, }\AttributeTok{lwd =} \DecValTok{2}\NormalTok{, }\AttributeTok{lty =} \DecValTok{2}\NormalTok{)}
\FunctionTok{legend}\NormalTok{(}\StringTok{"topright"}\NormalTok{, }\FunctionTok{c}\NormalTok{(}\FunctionTok{paste}\NormalTok{(}\StringTok{"Mean ="}\NormalTok{, }\FunctionTok{mean}\NormalTok{(x1)),}
                     \FunctionTok{paste}\NormalTok{(}\StringTok{"Median ="}\NormalTok{, }\FunctionTok{median}\NormalTok{(x1))),}
       \AttributeTok{col =} \FunctionTok{c}\NormalTok{(}\StringTok{"\#d7191c"}\NormalTok{, }\StringTok{"\#0571b0"}\NormalTok{), }\AttributeTok{lwd =} \DecValTok{2}\NormalTok{, }\AttributeTok{lty =} \FunctionTok{c}\NormalTok{(}\DecValTok{1}\NormalTok{, }\DecValTok{2}\NormalTok{), }\AttributeTok{bty =} \StringTok{"n"}\NormalTok{)}

\CommentTok{\# With outlier}
\NormalTok{x2 }\OtherTok{\textless{}{-}} \FunctionTok{c}\NormalTok{(x1, }\DecValTok{75}\NormalTok{)}
\FunctionTok{hist}\NormalTok{(x2, }\AttributeTok{col =} \StringTok{"\#f4a582"}\NormalTok{, }\AttributeTok{border =} \StringTok{"white"}\NormalTok{, }\AttributeTok{main =} \StringTok{"With Outlier (75)"}\NormalTok{,}
     \AttributeTok{xlab =} \StringTok{"Value"}\NormalTok{, }\AttributeTok{ylab =} \StringTok{""}\NormalTok{, }\AttributeTok{xlim =} \FunctionTok{c}\NormalTok{(}\DecValTok{0}\NormalTok{, }\DecValTok{80}\NormalTok{), }\AttributeTok{breaks =} \DecValTok{8}\NormalTok{)}
\FunctionTok{abline}\NormalTok{(}\AttributeTok{v =} \FunctionTok{mean}\NormalTok{(x2),   }\AttributeTok{col =} \StringTok{"\#d7191c"}\NormalTok{, }\AttributeTok{lwd =} \DecValTok{2}\NormalTok{, }\AttributeTok{lty =} \DecValTok{1}\NormalTok{)}
\FunctionTok{abline}\NormalTok{(}\AttributeTok{v =} \FunctionTok{median}\NormalTok{(x2), }\AttributeTok{col =} \StringTok{"\#0571b0"}\NormalTok{, }\AttributeTok{lwd =} \DecValTok{2}\NormalTok{, }\AttributeTok{lty =} \DecValTok{2}\NormalTok{)}
\FunctionTok{legend}\NormalTok{(}\StringTok{"topright"}\NormalTok{, }\FunctionTok{c}\NormalTok{(}\FunctionTok{paste}\NormalTok{(}\StringTok{"Mean ="}\NormalTok{, }\FunctionTok{round}\NormalTok{(}\FunctionTok{mean}\NormalTok{(x2), }\DecValTok{1}\NormalTok{)),}
                     \FunctionTok{paste}\NormalTok{(}\StringTok{"Median ="}\NormalTok{, }\FunctionTok{median}\NormalTok{(x2))),}
       \AttributeTok{col =} \FunctionTok{c}\NormalTok{(}\StringTok{"\#d7191c"}\NormalTok{, }\StringTok{"\#0571b0"}\NormalTok{), }\AttributeTok{lwd =} \DecValTok{2}\NormalTok{, }\AttributeTok{lty =} \FunctionTok{c}\NormalTok{(}\DecValTok{1}\NormalTok{, }\DecValTok{2}\NormalTok{), }\AttributeTok{bty =} \StringTok{"n"}\NormalTok{)}

\FunctionTok{par}\NormalTok{(}\AttributeTok{mfrow =} \FunctionTok{c}\NormalTok{(}\DecValTok{1}\NormalTok{, }\DecValTok{1}\NormalTok{))}
\end{Highlighting}
\end{Shaded}

\begin{figure}[H]

\centering{

\pandocbounded{\includegraphics[keepaspectratio]{chapters/04-measures_files/figure-latex/fig-mean-vs-median-1.pdf}}

}

\caption{\label{fig-mean-vs-median}Effect of outliers on mean vs
median.}

\end{figure}%

\section{Measures of spread}\label{measures-of-spread}

\subsection{Range}\label{range}

The \textbf{range} is the simplest measure of spread:

\[\text{Range} = \text{Maximum} - \text{Minimum}\]

\begin{Shaded}
\begin{Highlighting}[]
\FunctionTok{cat}\NormalTok{(}\StringTok{"Range of scores:"}\NormalTok{, }\FunctionTok{max}\NormalTok{(scores) }\SpecialCharTok{{-}} \FunctionTok{min}\NormalTok{(scores), }\StringTok{"}\SpecialCharTok{\textbackslash{}n}\StringTok{"}\NormalTok{)}
\CommentTok{\#\textgreater{} Range of scores: 43}
\FunctionTok{cat}\NormalTok{(}\StringTok{"  (Max ="}\NormalTok{, }\FunctionTok{max}\NormalTok{(scores), }\StringTok{", Min ="}\NormalTok{, }\FunctionTok{min}\NormalTok{(scores), }\StringTok{")}\SpecialCharTok{\textbackslash{}n}\StringTok{"}\NormalTok{)}
\CommentTok{\#\textgreater{}   (Max = 95 , Min = 52 )}
\end{Highlighting}
\end{Shaded}

\subsection{Quartiles and the interquartile range
(IQR)}\label{quartiles-and-the-interquartile-range-iqr}

\textbf{Quartiles} divide sorted data into four equal groups:

\begin{longtable}[]{@{}ll@{}}
\toprule\noalign{}
Quartile & Description \\
\midrule\noalign{}
\endhead
\bottomrule\noalign{}
\endlastfoot
\textbf{Q1} (lower quartile) & 25\% of data falls below this value \\
\textbf{Q2} (median) & 50\% of data falls below this value \\
\textbf{Q3} (upper quartile) & 75\% of data falls below this value \\
\end{longtable}

The \textbf{interquartile range (IQR)} is the spread of the middle 50\%
of data:

\[\text{IQR} = Q3 - Q1\]

\begin{Shaded}
\begin{Highlighting}[]
\NormalTok{q }\OtherTok{\textless{}{-}} \FunctionTok{quantile}\NormalTok{(scores)}
\FunctionTok{cat}\NormalTok{(}\StringTok{"Five{-}number summary:}\SpecialCharTok{\textbackslash{}n}\StringTok{"}\NormalTok{)}
\CommentTok{\#\textgreater{} Five{-}number summary:}
\FunctionTok{cat}\NormalTok{(}\StringTok{"  Minimum:"}\NormalTok{, q[}\DecValTok{1}\NormalTok{], }\StringTok{"}\SpecialCharTok{\textbackslash{}n}\StringTok{"}\NormalTok{)}
\CommentTok{\#\textgreater{}   Minimum: 52}
\FunctionTok{cat}\NormalTok{(}\StringTok{"  Q1:     "}\NormalTok{, q[}\DecValTok{2}\NormalTok{], }\StringTok{"}\SpecialCharTok{\textbackslash{}n}\StringTok{"}\NormalTok{)}
\CommentTok{\#\textgreater{}   Q1:      67.5}
\FunctionTok{cat}\NormalTok{(}\StringTok{"  Median: "}\NormalTok{, q[}\DecValTok{3}\NormalTok{], }\StringTok{"}\SpecialCharTok{\textbackslash{}n}\StringTok{"}\NormalTok{)}
\CommentTok{\#\textgreater{}   Median:  73.5}
\FunctionTok{cat}\NormalTok{(}\StringTok{"  Q3:     "}\NormalTok{, q[}\DecValTok{4}\NormalTok{], }\StringTok{"}\SpecialCharTok{\textbackslash{}n}\StringTok{"}\NormalTok{)}
\CommentTok{\#\textgreater{}   Q3:      80.75}
\FunctionTok{cat}\NormalTok{(}\StringTok{"  Maximum:"}\NormalTok{, q[}\DecValTok{5}\NormalTok{], }\StringTok{"}\SpecialCharTok{\textbackslash{}n}\StringTok{"}\NormalTok{)}
\CommentTok{\#\textgreater{}   Maximum: 95}
\FunctionTok{cat}\NormalTok{(}\StringTok{"  IQR ="}\NormalTok{, q[}\DecValTok{4}\NormalTok{] }\SpecialCharTok{{-}}\NormalTok{ q[}\DecValTok{2}\NormalTok{], }\StringTok{"}\SpecialCharTok{\textbackslash{}n}\StringTok{"}\NormalTok{)}
\CommentTok{\#\textgreater{}   IQR = 13.25}
\end{Highlighting}
\end{Shaded}

\subsection{Standard deviation}\label{standard-deviation}

The \textbf{standard deviation} measures how spread out data values are
from the mean. A larger standard deviation means more variability.

\[s = \sqrt{\frac{\sum (x_i - \bar{x})^2}{n - 1}}\]

\begin{Shaded}
\begin{Highlighting}[]
\FunctionTok{cat}\NormalTok{(}\StringTok{"Mean:"}\NormalTok{, }\FunctionTok{round}\NormalTok{(}\FunctionTok{mean}\NormalTok{(scores), }\DecValTok{2}\NormalTok{), }\StringTok{"}\SpecialCharTok{\textbackslash{}n}\StringTok{"}\NormalTok{)}
\CommentTok{\#\textgreater{} Mean: 73.85}
\FunctionTok{cat}\NormalTok{(}\StringTok{"Standard deviation:"}\NormalTok{, }\FunctionTok{round}\NormalTok{(}\FunctionTok{sd}\NormalTok{(scores), }\DecValTok{2}\NormalTok{), }\StringTok{"}\SpecialCharTok{\textbackslash{}n}\StringTok{"}\NormalTok{)}
\CommentTok{\#\textgreater{} Standard deviation: 11.52}
\end{Highlighting}
\end{Shaded}

\begin{tcolorbox}[enhanced jigsaw, bottomtitle=1mm, opacityback=0, titlerule=0mm, leftrule=.75mm, opacitybacktitle=0.6, coltitle=black, toprule=.15mm, colbacktitle=quarto-callout-tip-color!10!white, title=\textcolor{quarto-callout-tip-color}{\faLightbulb}\hspace{0.5em}{Rule of thumb}, left=2mm, rightrule=.15mm, toptitle=1mm, colframe=quarto-callout-tip-color-frame, breakable, arc=.35mm, bottomrule=.15mm, colback=white]

For roughly symmetric distributions, about \textbf{68\%} of data falls
within one standard deviation of the mean, and \textbf{95\%} falls
within two standard deviations (the empirical rule / 68-95-99.7 rule).

\end{tcolorbox}

\subsection{Choosing the right measure of
spread}\label{choosing-the-right-measure-of-spread}

\begin{longtable}[]{@{}
  >{\raggedright\arraybackslash}p{(\linewidth - 4\tabcolsep) * \real{0.2812}}
  >{\raggedright\arraybackslash}p{(\linewidth - 4\tabcolsep) * \real{0.3750}}
  >{\raggedright\arraybackslash}p{(\linewidth - 4\tabcolsep) * \real{0.3438}}@{}}
\toprule\noalign{}
\begin{minipage}[b]{\linewidth}\raggedright
Measure
\end{minipage} & \begin{minipage}[b]{\linewidth}\raggedright
Use with\ldots{}
\end{minipage} & \begin{minipage}[b]{\linewidth}\raggedright
Advantage
\end{minipage} \\
\midrule\noalign{}
\endhead
\bottomrule\noalign{}
\endlastfoot
\textbf{Range} & Quick overview & Easy to calculate \\
\textbf{IQR} & Skewed data or when outliers present & Resistant to
outliers \\
\textbf{Standard deviation} & Symmetric data, no outliers & Uses all
data values \\
\end{longtable}

\section{Comparing groups}\label{comparing-groups}

A common task in AS91946 is to \textbf{compare two or more groups}.
Always compare both centre AND spread.

\begin{Shaded}
\begin{Highlighting}[]
\FunctionTok{set.seed}\NormalTok{(}\DecValTok{42}\NormalTok{)}
\NormalTok{group\_a }\OtherTok{\textless{}{-}} \FunctionTok{round}\NormalTok{(}\FunctionTok{rnorm}\NormalTok{(}\DecValTok{30}\NormalTok{, }\AttributeTok{mean =} \DecValTok{65}\NormalTok{, }\AttributeTok{sd =} \DecValTok{10}\NormalTok{))}
\NormalTok{group\_b }\OtherTok{\textless{}{-}} \FunctionTok{round}\NormalTok{(}\FunctionTok{rnorm}\NormalTok{(}\DecValTok{30}\NormalTok{, }\AttributeTok{mean =} \DecValTok{72}\NormalTok{, }\AttributeTok{sd =} \DecValTok{15}\NormalTok{))}
\end{Highlighting}
\end{Shaded}

\begin{Shaded}
\begin{Highlighting}[]
\FunctionTok{boxplot}\NormalTok{(}
  \FunctionTok{list}\NormalTok{(}\StringTok{"Group A"} \OtherTok{=}\NormalTok{ group\_a, }\StringTok{"Group B"} \OtherTok{=}\NormalTok{ group\_b),}
  \AttributeTok{col    =} \FunctionTok{c}\NormalTok{(}\StringTok{"\#4dac26"}\NormalTok{, }\StringTok{"\#0571b0"}\NormalTok{),}
  \AttributeTok{border =} \StringTok{"\#333333"}\NormalTok{,}
  \AttributeTok{main   =} \StringTok{"Maths Scores: Group A vs Group B"}\NormalTok{,}
  \AttributeTok{ylab   =} \StringTok{"Score"}\NormalTok{,}
  \AttributeTok{las    =} \DecValTok{1}
\NormalTok{)}
\end{Highlighting}
\end{Shaded}

\begin{figure}[H]

\centering{

\pandocbounded{\includegraphics[keepaspectratio]{chapters/04-measures_files/figure-latex/fig-compare-boxplot-1.pdf}}

}

\caption{\label{fig-compare-boxplot}Comparing maths scores between Group
A and Group B.}

\end{figure}%

\begin{Shaded}
\begin{Highlighting}[]
\NormalTok{summary\_stats }\OtherTok{\textless{}{-}} \FunctionTok{data.frame}\NormalTok{(}
  \AttributeTok{Group    =} \FunctionTok{c}\NormalTok{(}\StringTok{"Group A"}\NormalTok{, }\StringTok{"Group B"}\NormalTok{),}
  \AttributeTok{n        =} \FunctionTok{c}\NormalTok{(}\FunctionTok{length}\NormalTok{(group\_a), }\FunctionTok{length}\NormalTok{(group\_b)),}
  \AttributeTok{Mean     =} \FunctionTok{round}\NormalTok{(}\FunctionTok{c}\NormalTok{(}\FunctionTok{mean}\NormalTok{(group\_a), }\FunctionTok{mean}\NormalTok{(group\_b)), }\DecValTok{1}\NormalTok{),}
  \AttributeTok{Median   =} \FunctionTok{c}\NormalTok{(}\FunctionTok{median}\NormalTok{(group\_a), }\FunctionTok{median}\NormalTok{(group\_b)),}
  \AttributeTok{SD       =} \FunctionTok{round}\NormalTok{(}\FunctionTok{c}\NormalTok{(}\FunctionTok{sd}\NormalTok{(group\_a), }\FunctionTok{sd}\NormalTok{(group\_b)), }\DecValTok{1}\NormalTok{),}
  \AttributeTok{IQR      =} \FunctionTok{c}\NormalTok{(}\FunctionTok{IQR}\NormalTok{(group\_a), }\FunctionTok{IQR}\NormalTok{(group\_b)),}
  \AttributeTok{Min      =} \FunctionTok{c}\NormalTok{(}\FunctionTok{min}\NormalTok{(group\_a), }\FunctionTok{min}\NormalTok{(group\_b)),}
  \AttributeTok{Max      =} \FunctionTok{c}\NormalTok{(}\FunctionTok{max}\NormalTok{(group\_a), }\FunctionTok{max}\NormalTok{(group\_b))}
\NormalTok{)}

\NormalTok{knitr}\SpecialCharTok{::}\FunctionTok{kable}\NormalTok{(summary\_stats, }\AttributeTok{caption =} \StringTok{"Summary statistics for Groups A and B"}\NormalTok{)}
\end{Highlighting}
\end{Shaded}

\begin{longtable}[]{@{}lrrrrrrr@{}}
\caption{Summary statistics for Groups A and B}\tabularnewline
\toprule\noalign{}
Group & n & Mean & Median & SD & IQR & Min & Max \\
\midrule\noalign{}
\endfirsthead
\toprule\noalign{}
Group & n & Mean & Median & SD & IQR & Min & Max \\
\midrule\noalign{}
\endhead
\bottomrule\noalign{}
\endlastfoot
Group A & 30 & 65.6 & 64 & 12.6 & 14.25 & 38 & 88 \\
Group B & 30 & 70.2 & 74 & 15.8 & 21.25 & 27 & 96 \\
\end{longtable}

\subsection{Model comparison paragraph (Merit
level)}\label{model-comparison-paragraph-merit-level}

\begin{quote}
``Group B has a higher median score (72) than Group A (65), suggesting
that Group B performed better overall. However, Group B also has a
larger IQR (approximately 20 vs 14) and standard deviation (15.1 vs
9.8), indicating their scores were more variable. Group A's performance
was more consistent. Both groups appear to have roughly symmetric
distributions with no obvious outliers.''
\end{quote}

\section{Identifying outliers}\label{identifying-outliers}

A value is often flagged as an \textbf{outlier} if it falls more than
1.5 × IQR below Q1 or above Q3:

\[\text{Lower fence} = Q1 - 1.5 \times \text{IQR}\]
\[\text{Upper fence} = Q3 + 1.5 \times \text{IQR}\]

\begin{Shaded}
\begin{Highlighting}[]
\NormalTok{Q1  }\OtherTok{\textless{}{-}} \FunctionTok{quantile}\NormalTok{(scores, }\FloatTok{0.25}\NormalTok{)}
\NormalTok{Q3  }\OtherTok{\textless{}{-}} \FunctionTok{quantile}\NormalTok{(scores, }\FloatTok{0.75}\NormalTok{)}
\NormalTok{iqr }\OtherTok{\textless{}{-}}\NormalTok{ Q3 }\SpecialCharTok{{-}}\NormalTok{ Q1}

\NormalTok{lower\_fence }\OtherTok{\textless{}{-}}\NormalTok{ Q1 }\SpecialCharTok{{-}} \FloatTok{1.5} \SpecialCharTok{*}\NormalTok{ iqr}
\NormalTok{upper\_fence }\OtherTok{\textless{}{-}}\NormalTok{ Q3 }\SpecialCharTok{+} \FloatTok{1.5} \SpecialCharTok{*}\NormalTok{ iqr}

\FunctionTok{cat}\NormalTok{(}\StringTok{"Q1 ="}\NormalTok{, Q1, }\StringTok{"}\SpecialCharTok{\textbackslash{}n}\StringTok{"}\NormalTok{)}
\CommentTok{\#\textgreater{} Q1 = 67.5}
\FunctionTok{cat}\NormalTok{(}\StringTok{"Q3 ="}\NormalTok{, Q3, }\StringTok{"}\SpecialCharTok{\textbackslash{}n}\StringTok{"}\NormalTok{)}
\CommentTok{\#\textgreater{} Q3 = 80.75}
\FunctionTok{cat}\NormalTok{(}\StringTok{"IQR ="}\NormalTok{, iqr, }\StringTok{"}\SpecialCharTok{\textbackslash{}n}\StringTok{"}\NormalTok{)}
\CommentTok{\#\textgreater{} IQR = 13.25}
\FunctionTok{cat}\NormalTok{(}\StringTok{"Lower fence:"}\NormalTok{, lower\_fence, }\StringTok{"}\SpecialCharTok{\textbackslash{}n}\StringTok{"}\NormalTok{)}
\CommentTok{\#\textgreater{} Lower fence: 47.625}
\FunctionTok{cat}\NormalTok{(}\StringTok{"Upper fence:"}\NormalTok{, upper\_fence, }\StringTok{"}\SpecialCharTok{\textbackslash{}n}\StringTok{"}\NormalTok{)}
\CommentTok{\#\textgreater{} Upper fence: 100.625}
\FunctionTok{cat}\NormalTok{(}\StringTok{"Outliers:"}\NormalTok{, scores[scores }\SpecialCharTok{\textless{}}\NormalTok{ lower\_fence }\SpecialCharTok{|}\NormalTok{ scores }\SpecialCharTok{\textgreater{}}\NormalTok{ upper\_fence], }\StringTok{"}\SpecialCharTok{\textbackslash{}n}\StringTok{"}\NormalTok{)}
\CommentTok{\#\textgreater{} Outliers:}
\end{Highlighting}
\end{Shaded}

\section{Glossary}\label{glossary-2}

\begin{longtable}[]{@{}ll@{}}
\toprule\noalign{}
Term & Meaning \\
\midrule\noalign{}
\endhead
\bottomrule\noalign{}
\endlastfoot
\textbf{Mean} & Sum of all values divided by the count \\
\textbf{Median} & Middle value when data is sorted \\
\textbf{Mode} & Most frequently occurring value \\
\textbf{Range} & Maximum − Minimum \\
\textbf{IQR} & Q3 − Q1; spread of the middle 50\% \\
\textbf{Standard deviation} & Average distance of values from the
mean \\
\textbf{Outlier} & A value that lies far from the rest of the data \\
\textbf{Five-number summary} & Min, Q1, Median, Q3, Max \\
\end{longtable}

\section{Check your understanding}\label{check-your-understanding-3}

Use the dataset: \textbf{12, 15, 18, 18, 20, 22, 25, 25, 28, 100}

\begin{enumerate}
\def\labelenumi{\arabic{enumi}.}
\tightlist
\item
  Calculate the mean and median. Which is more appropriate as a measure
  of centre for this data? Why?
\item
  Find Q1, Q3, and the IQR.
\item
  Calculate the range.
\item
  Is 100 an outlier? Use the 1.5 × IQR rule to justify your answer.
\item
  Sketch a box plot of this data.
\end{enumerate}

\begin{tcolorbox}[enhanced jigsaw, bottomtitle=1mm, opacityback=0, titlerule=0mm, leftrule=.75mm, opacitybacktitle=0.6, coltitle=black, toprule=.15mm, colbacktitle=quarto-callout-note-color!10!white, title=\textcolor{quarto-callout-note-color}{\faInfo}\hspace{0.5em}{Answers}, left=2mm, rightrule=.15mm, toptitle=1mm, colframe=quarto-callout-note-color-frame, breakable, arc=.35mm, bottomrule=.15mm, colback=white]

\begin{enumerate}
\def\labelenumi{\arabic{enumi}.}
\tightlist
\item
  Mean = (12+15+18+18+20+22+25+25+28+100)/10 = 283/10 = 28.3; Median =
  (20+22)/2 = 21. The \textbf{median} is more appropriate because of the
  outlier (100) which pulls the mean up.
\item
  Q1 = 18, Q3 = 25, IQR = 7.
\item
  Range = 100 − 12 = 88.
\item
  Upper fence = 25 + 1.5 × 7 = 25 + 10.5 = 35.5. Since 100
  \textgreater{} 35.5, \textbf{yes}, 100 is an outlier.
\item
  Sketch: whiskers from 12 to 35.5, box from 18 to 25, median line at
  21, and a dot at 100 representing the outlier.
\end{enumerate}

\end{tcolorbox}

\bookmarksetup{startatroot}

\chapter{Probability}\label{probability}

\bookmarksetup{startatroot}

\chapter{Probability}\label{sec-probability}

\textbf{Probability} is a measure of how likely an event is to occur. It
is expressed as a number between 0 (impossible) and 1 (certain), or as a
percentage.

\[P(\text{event}) = \frac{\text{number of favourable outcomes}}{\text{total number of possible outcomes}}\]

\section{Probability scale}\label{probability-scale}

\begin{figure}

\centering{

\pandocbounded{\includegraphics[keepaspectratio]{chapters/05-probability_files/figure-latex/fig-prob-scale-1.pdf}}

}

\caption{\label{fig-prob-scale}The probability scale from impossible (0)
to certain (1).}

\end{figure}%

\section{Key vocabulary}\label{key-vocabulary}

\begin{longtable}[]{@{}ll@{}}
\toprule\noalign{}
Term & Definition \\
\midrule\noalign{}
\endhead
\bottomrule\noalign{}
\endlastfoot
\textbf{Experiment} & A process with uncertain outcomes \\
\textbf{Outcome} & A single possible result \\
\textbf{Event} & One or more outcomes grouped together \\
\textbf{Sample space} & The set of all possible outcomes \\
\textbf{Complementary event} & All outcomes that are NOT in the event \\
\end{longtable}

\section{Simple probability --- equally likely
outcomes}\label{simple-probability-equally-likely-outcomes}

When all outcomes are \textbf{equally likely}, we can use the classical
definition.

\subsection{Example: rolling a die}\label{example-rolling-a-die}

\begin{Shaded}
\begin{Highlighting}[]
\NormalTok{sample\_space }\OtherTok{\textless{}{-}} \DecValTok{1}\SpecialCharTok{:}\DecValTok{6}
\FunctionTok{cat}\NormalTok{(}\StringTok{"Sample space:"}\NormalTok{, sample\_space, }\StringTok{"}\SpecialCharTok{\textbackslash{}n}\StringTok{"}\NormalTok{)}
\CommentTok{\#\textgreater{} Sample space: 1 2 3 4 5 6}

\CommentTok{\# P(rolling a 4)}
\NormalTok{p\_four }\OtherTok{\textless{}{-}} \FunctionTok{sum}\NormalTok{(sample\_space }\SpecialCharTok{==} \DecValTok{4}\NormalTok{) }\SpecialCharTok{/} \FunctionTok{length}\NormalTok{(sample\_space)}
\FunctionTok{cat}\NormalTok{(}\StringTok{"P(rolling a 4) ="}\NormalTok{, p\_four, }\StringTok{"}\SpecialCharTok{\textbackslash{}n}\StringTok{"}\NormalTok{)}
\CommentTok{\#\textgreater{} P(rolling a 4) = 0.1666667}

\CommentTok{\# P(rolling an even number)}
\NormalTok{p\_even }\OtherTok{\textless{}{-}} \FunctionTok{sum}\NormalTok{(sample\_space }\SpecialCharTok{\%\%} \DecValTok{2} \SpecialCharTok{==} \DecValTok{0}\NormalTok{) }\SpecialCharTok{/} \FunctionTok{length}\NormalTok{(sample\_space)}
\FunctionTok{cat}\NormalTok{(}\StringTok{"P(even number) ="}\NormalTok{, p\_even, }\StringTok{"}\SpecialCharTok{\textbackslash{}n}\StringTok{"}\NormalTok{)}
\CommentTok{\#\textgreater{} P(even number) = 0.5}

\CommentTok{\# P(rolling \textgreater{} 4)}
\NormalTok{p\_gt4 }\OtherTok{\textless{}{-}} \FunctionTok{sum}\NormalTok{(sample\_space }\SpecialCharTok{\textgreater{}} \DecValTok{4}\NormalTok{) }\SpecialCharTok{/} \FunctionTok{length}\NormalTok{(sample\_space)}
\FunctionTok{cat}\NormalTok{(}\StringTok{"P(\textgreater{} 4) ="}\NormalTok{, p\_gt4, }\StringTok{"}\SpecialCharTok{\textbackslash{}n}\StringTok{"}\NormalTok{)}
\CommentTok{\#\textgreater{} P(\textgreater{} 4) = 0.3333333}
\end{Highlighting}
\end{Shaded}

\section{Complementary events}\label{complementary-events}

The \textbf{complement} of event A (written A' or \(\bar{A}\)) contains
all outcomes NOT in A.

\[P(A') = 1 - P(A)\]

\textbf{Example:} If the probability of rain tomorrow is 0.3, then:

\[P(\text{no rain}) = 1 - 0.3 = 0.7\]

\begin{Shaded}
\begin{Highlighting}[]
\NormalTok{p\_rain }\OtherTok{\textless{}{-}} \FloatTok{0.3}
\NormalTok{p\_no\_rain }\OtherTok{\textless{}{-}} \DecValTok{1} \SpecialCharTok{{-}}\NormalTok{ p\_rain}
\FunctionTok{cat}\NormalTok{(}\StringTok{"P(rain) ="}\NormalTok{, p\_rain, }\StringTok{"}\SpecialCharTok{\textbackslash{}n}\StringTok{"}\NormalTok{)}
\CommentTok{\#\textgreater{} P(rain) = 0.3}
\FunctionTok{cat}\NormalTok{(}\StringTok{"P(no rain) = 1 {-}"}\NormalTok{, p\_rain, }\StringTok{"="}\NormalTok{, p\_no\_rain, }\StringTok{"}\SpecialCharTok{\textbackslash{}n}\StringTok{"}\NormalTok{)}
\CommentTok{\#\textgreater{} P(no rain) = 1 {-} 0.3 = 0.7}
\end{Highlighting}
\end{Shaded}

\section{Experimental (relative frequency)
probability}\label{experimental-relative-frequency-probability}

When we cannot calculate probability theoretically, we \textbf{estimate}
it by repeating an experiment many times:

\[P(\text{event}) \approx \frac{\text{number of times event occurred}}{\text{total number of trials}}\]

\subsection{Simulating coin flips}\label{simulating-coin-flips}

\begin{Shaded}
\begin{Highlighting}[]
\FunctionTok{set.seed}\NormalTok{(}\DecValTok{99}\NormalTok{)}
\NormalTok{flips  }\OtherTok{\textless{}{-}} \FunctionTok{sample}\NormalTok{(}\FunctionTok{c}\NormalTok{(}\StringTok{"Heads"}\NormalTok{, }\StringTok{"Tails"}\NormalTok{), }\DecValTok{1000}\NormalTok{, }\AttributeTok{replace =} \ConstantTok{TRUE}\NormalTok{)}
\NormalTok{cumulative\_heads }\OtherTok{\textless{}{-}} \FunctionTok{cumsum}\NormalTok{(flips }\SpecialCharTok{==} \StringTok{"Heads"}\NormalTok{) }\SpecialCharTok{/} \FunctionTok{seq\_along}\NormalTok{(flips)}

\FunctionTok{plot}\NormalTok{(}
\NormalTok{  cumulative\_heads,}
  \AttributeTok{type =} \StringTok{"l"}\NormalTok{,}
  \AttributeTok{col  =} \StringTok{"\#0571b0"}\NormalTok{,}
  \AttributeTok{lwd  =} \FloatTok{1.5}\NormalTok{,}
  \AttributeTok{xlab =} \StringTok{"Number of Flips"}\NormalTok{,}
  \AttributeTok{ylab =} \StringTok{"Relative Frequency of Heads"}\NormalTok{,}
  \AttributeTok{main =} \StringTok{"Simulated Coin Flips"}\NormalTok{,}
  \AttributeTok{ylim =} \FunctionTok{c}\NormalTok{(}\DecValTok{0}\NormalTok{, }\DecValTok{1}\NormalTok{),}
  \AttributeTok{las  =} \DecValTok{1}
\NormalTok{)}
\FunctionTok{abline}\NormalTok{(}\AttributeTok{h =} \FloatTok{0.5}\NormalTok{, }\AttributeTok{col =} \StringTok{"\#d7191c"}\NormalTok{, }\AttributeTok{lwd =} \DecValTok{2}\NormalTok{, }\AttributeTok{lty =} \DecValTok{2}\NormalTok{)}
\FunctionTok{legend}\NormalTok{(}\StringTok{"topright"}\NormalTok{, }\FunctionTok{c}\NormalTok{(}\StringTok{"Relative frequency"}\NormalTok{, }\StringTok{"True probability (0.5)"}\NormalTok{),}
       \AttributeTok{col =} \FunctionTok{c}\NormalTok{(}\StringTok{"\#0571b0"}\NormalTok{, }\StringTok{"\#d7191c"}\NormalTok{), }\AttributeTok{lwd =} \DecValTok{2}\NormalTok{, }\AttributeTok{lty =} \FunctionTok{c}\NormalTok{(}\DecValTok{1}\NormalTok{, }\DecValTok{2}\NormalTok{), }\AttributeTok{bty =} \StringTok{"n"}\NormalTok{)}
\end{Highlighting}
\end{Shaded}

\begin{figure}[H]

\centering{

\pandocbounded{\includegraphics[keepaspectratio]{chapters/05-probability_files/figure-latex/fig-coin-simulation-1.pdf}}

}

\caption{\label{fig-coin-simulation}Relative frequency of heads
approaches 0.5 as the number of flips increases (the law of large
numbers).}

\end{figure}%

\begin{tcolorbox}[enhanced jigsaw, bottomtitle=1mm, opacityback=0, titlerule=0mm, leftrule=.75mm, opacitybacktitle=0.6, coltitle=black, toprule=.15mm, colbacktitle=quarto-callout-note-color!10!white, title=\textcolor{quarto-callout-note-color}{\faInfo}\hspace{0.5em}{Law of Large Numbers}, left=2mm, rightrule=.15mm, toptitle=1mm, colframe=quarto-callout-note-color-frame, breakable, arc=.35mm, bottomrule=.15mm, colback=white]

As the number of trials increases, the \textbf{relative frequency}
(experimental probability) gets closer and closer to the \textbf{true
probability}. This is why surveys and experiments need large samples.

\end{tcolorbox}

\section{Two-way tables}\label{two-way-tables}

A \textbf{two-way table} (contingency table) organises data by two
categorical variables and is very useful for calculating probabilities.

\begin{Shaded}
\begin{Highlighting}[]
\CommentTok{\# Simulated data: 120 students, sport preference by gender}
\FunctionTok{set.seed}\NormalTok{(}\DecValTok{7}\NormalTok{)}
\NormalTok{n\_students }\OtherTok{\textless{}{-}} \DecValTok{120}
\NormalTok{sport\_data }\OtherTok{\textless{}{-}} \FunctionTok{data.frame}\NormalTok{(}
  \AttributeTok{Gender =} \FunctionTok{sample}\NormalTok{(}\FunctionTok{c}\NormalTok{(}\StringTok{"Male"}\NormalTok{, }\StringTok{"Female"}\NormalTok{), n\_students, }\AttributeTok{replace =} \ConstantTok{TRUE}\NormalTok{,}
                  \AttributeTok{prob =} \FunctionTok{c}\NormalTok{(}\FloatTok{0.5}\NormalTok{, }\FloatTok{0.5}\NormalTok{)),}
  \AttributeTok{Sport  =} \FunctionTok{sample}\NormalTok{(}\FunctionTok{c}\NormalTok{(}\StringTok{"Rugby"}\NormalTok{, }\StringTok{"Netball"}\NormalTok{, }\StringTok{"Football"}\NormalTok{, }\StringTok{"Basketball"}\NormalTok{),}
\NormalTok{                  n\_students, }\AttributeTok{replace =} \ConstantTok{TRUE}\NormalTok{, }\AttributeTok{prob =} \FunctionTok{c}\NormalTok{(}\FloatTok{0.3}\NormalTok{, }\FloatTok{0.25}\NormalTok{, }\FloatTok{0.25}\NormalTok{, }\FloatTok{0.2}\NormalTok{))}
\NormalTok{)}

\NormalTok{tbl }\OtherTok{\textless{}{-}} \FunctionTok{table}\NormalTok{(sport\_data}\SpecialCharTok{$}\NormalTok{Sport, sport\_data}\SpecialCharTok{$}\NormalTok{Gender)}
\NormalTok{tbl\_with\_total }\OtherTok{\textless{}{-}} \FunctionTok{addmargins}\NormalTok{(tbl)}
\NormalTok{knitr}\SpecialCharTok{::}\FunctionTok{kable}\NormalTok{(tbl\_with\_total, }\AttributeTok{caption =} \StringTok{"Sport preference by gender (n = 120)"}\NormalTok{)}
\end{Highlighting}
\end{Shaded}

\begin{longtable}[]{@{}lrrr@{}}
\caption{Sport preference by gender (n = 120)}\tabularnewline
\toprule\noalign{}
& Female & Male & Sum \\
\midrule\noalign{}
\endfirsthead
\toprule\noalign{}
& Female & Male & Sum \\
\midrule\noalign{}
\endhead
\bottomrule\noalign{}
\endlastfoot
Basketball & 16 & 11 & 27 \\
Football & 13 & 17 & 30 \\
Netball & 18 & 17 & 35 \\
Rugby & 16 & 12 & 28 \\
Sum & 63 & 57 & 120 \\
\end{longtable}

\subsection{Reading probabilities from a two-way
table}\label{reading-probabilities-from-a-two-way-table}

Using the table above:

\begin{Shaded}
\begin{Highlighting}[]
\NormalTok{total }\OtherTok{\textless{}{-}} \FunctionTok{sum}\NormalTok{(tbl)}

\FunctionTok{cat}\NormalTok{(}\StringTok{"P(Male) ="}\NormalTok{, }\FunctionTok{sum}\NormalTok{(tbl[, }\StringTok{"Male"}\NormalTok{]), }\StringTok{"/"}\NormalTok{, total, }\StringTok{"="}\NormalTok{,}
    \FunctionTok{round}\NormalTok{(}\FunctionTok{sum}\NormalTok{(tbl[, }\StringTok{"Male"}\NormalTok{]) }\SpecialCharTok{/}\NormalTok{ total, }\DecValTok{3}\NormalTok{), }\StringTok{"}\SpecialCharTok{\textbackslash{}n}\StringTok{"}\NormalTok{)}
\CommentTok{\#\textgreater{} P(Male) = 57 / 120 = 0.475}

\FunctionTok{cat}\NormalTok{(}\StringTok{"P(Rugby) ="}\NormalTok{, }\FunctionTok{sum}\NormalTok{(tbl[}\StringTok{"Rugby"}\NormalTok{, ]), }\StringTok{"/"}\NormalTok{, total, }\StringTok{"="}\NormalTok{,}
    \FunctionTok{round}\NormalTok{(}\FunctionTok{sum}\NormalTok{(tbl[}\StringTok{"Rugby"}\NormalTok{, ]) }\SpecialCharTok{/}\NormalTok{ total, }\DecValTok{3}\NormalTok{), }\StringTok{"}\SpecialCharTok{\textbackslash{}n}\StringTok{"}\NormalTok{)}
\CommentTok{\#\textgreater{} P(Rugby) = 28 / 120 = 0.233}

\FunctionTok{cat}\NormalTok{(}\StringTok{"P(Male AND Rugby) ="}\NormalTok{, tbl[}\StringTok{"Rugby"}\NormalTok{, }\StringTok{"Male"}\NormalTok{], }\StringTok{"/"}\NormalTok{, total, }\StringTok{"="}\NormalTok{,}
    \FunctionTok{round}\NormalTok{(tbl[}\StringTok{"Rugby"}\NormalTok{, }\StringTok{"Male"}\NormalTok{] }\SpecialCharTok{/}\NormalTok{ total, }\DecValTok{3}\NormalTok{), }\StringTok{"}\SpecialCharTok{\textbackslash{}n}\StringTok{"}\NormalTok{)}
\CommentTok{\#\textgreater{} P(Male AND Rugby) = 12 / 120 = 0.1}

\FunctionTok{cat}\NormalTok{(}\StringTok{"P(Rugby | Male) ="}\NormalTok{, tbl[}\StringTok{"Rugby"}\NormalTok{, }\StringTok{"Male"}\NormalTok{], }\StringTok{"/"}\NormalTok{,}
    \FunctionTok{sum}\NormalTok{(tbl[, }\StringTok{"Male"}\NormalTok{]), }\StringTok{"="}\NormalTok{,}
    \FunctionTok{round}\NormalTok{(tbl[}\StringTok{"Rugby"}\NormalTok{, }\StringTok{"Male"}\NormalTok{] }\SpecialCharTok{/} \FunctionTok{sum}\NormalTok{(tbl[, }\StringTok{"Male"}\NormalTok{]), }\DecValTok{3}\NormalTok{), }\StringTok{"}\SpecialCharTok{\textbackslash{}n}\StringTok{"}\NormalTok{)}
\CommentTok{\#\textgreater{} P(Rugby | Male) = 12 / 57 = 0.211}
\end{Highlighting}
\end{Shaded}

The last probability --- P(Rugby \textbar{} Male) --- is a
\textbf{conditional probability}: the probability of preferring rugby,
\textbf{given} the student is male.

\section{Tree diagrams}\label{tree-diagrams}

Tree diagrams help calculate probabilities for \textbf{sequential
events}.

\begin{figure}

\centering{

\pandocbounded{\includegraphics[keepaspectratio]{chapters/05-probability_files/figure-latex/fig-tree-diagram-1.pdf}}

}

\caption{\label{fig-tree-diagram}Tree diagram for drawing two balls from
a bag (with replacement).}

\end{figure}%

\textbf{Key rule:} Multiply probabilities along each branch, add
probabilities for different paths leading to the same outcome.

\section{Expected value (extension)}\label{expected-value-extension}

For a simple scenario, the \textbf{expected value} is the long-run
average outcome:

\[E(X) = \sum x \cdot P(x)\]

\textbf{Example:} A raffle has 200 tickets at \$2 each. First prize is
\$100, second prize is \$50. What is the expected gain per ticket?

\begin{Shaded}
\begin{Highlighting}[]
\NormalTok{outcomes  }\OtherTok{\textless{}{-}} \FunctionTok{c}\NormalTok{(}\DecValTok{100}\NormalTok{, }\DecValTok{50}\NormalTok{, }\DecValTok{0}\NormalTok{)         }\CommentTok{\# winnings}
\NormalTok{probs     }\OtherTok{\textless{}{-}} \FunctionTok{c}\NormalTok{(}\DecValTok{1}\SpecialCharTok{/}\DecValTok{200}\NormalTok{, }\DecValTok{1}\SpecialCharTok{/}\DecValTok{200}\NormalTok{, }\DecValTok{198}\SpecialCharTok{/}\DecValTok{200}\NormalTok{)}
\NormalTok{cost      }\OtherTok{\textless{}{-}} \DecValTok{2}

\NormalTok{expected\_win  }\OtherTok{\textless{}{-}} \FunctionTok{sum}\NormalTok{(outcomes }\SpecialCharTok{*}\NormalTok{ probs)}
\NormalTok{expected\_gain }\OtherTok{\textless{}{-}}\NormalTok{ expected\_win }\SpecialCharTok{{-}}\NormalTok{ cost}

\FunctionTok{cat}\NormalTok{(}\StringTok{"Expected winnings:"}\NormalTok{, expected\_win, }\StringTok{"}\SpecialCharTok{\textbackslash{}n}\StringTok{"}\NormalTok{)}
\CommentTok{\#\textgreater{} Expected winnings: 0.75}
\FunctionTok{cat}\NormalTok{(}\StringTok{"Ticket cost:      "}\NormalTok{, cost, }\StringTok{"}\SpecialCharTok{\textbackslash{}n}\StringTok{"}\NormalTok{)}
\CommentTok{\#\textgreater{} Ticket cost:       2}
\FunctionTok{cat}\NormalTok{(}\StringTok{"Expected net gain:"}\NormalTok{, expected\_gain, }\StringTok{"}\SpecialCharTok{\textbackslash{}n}\StringTok{"}\NormalTok{)}
\CommentTok{\#\textgreater{} Expected net gain: {-}1.25}
\FunctionTok{cat}\NormalTok{(}\StringTok{"(A negative expected gain means this is not a good investment!)}\SpecialCharTok{\textbackslash{}n}\StringTok{"}\NormalTok{)}
\CommentTok{\#\textgreater{} (A negative expected gain means this is not a good investment!)}
\end{Highlighting}
\end{Shaded}

\section{Glossary}\label{glossary-3}

\begin{longtable}[]{@{}
  >{\raggedright\arraybackslash}p{(\linewidth - 2\tabcolsep) * \real{0.3333}}
  >{\raggedright\arraybackslash}p{(\linewidth - 2\tabcolsep) * \real{0.6667}}@{}}
\toprule\noalign{}
\begin{minipage}[b]{\linewidth}\raggedright
Term
\end{minipage} & \begin{minipage}[b]{\linewidth}\raggedright
Definition
\end{minipage} \\
\midrule\noalign{}
\endhead
\bottomrule\noalign{}
\endlastfoot
\textbf{Probability} & Likelihood of an event, from 0 to 1 \\
\textbf{Sample space} & All possible outcomes \\
\textbf{Event} & One or more outcomes \\
\textbf{Complement} & All outcomes NOT in the event \\
\textbf{Relative frequency} & Experimental estimate of probability \\
\textbf{Conditional probability} & P(A given B): probability of A given
B occurred \\
\textbf{Two-way table} & Table showing two categorical variables
together \\
\textbf{Tree diagram} & Diagram showing sequential outcomes and
probabilities \\
\end{longtable}

\section{Check your understanding}\label{check-your-understanding-4}

\begin{enumerate}
\def\labelenumi{\arabic{enumi}.}
\tightlist
\item
  A bag contains 4 red, 3 blue, and 3 yellow marbles. A marble is drawn
  at random. Find:

  \begin{enumerate}
  \def\labelenumii{\alph{enumii}.}
  \tightlist
  \item
    P(red)
  \item
    P(not red)
  \item
    P(blue or yellow)
  \end{enumerate}
\item
  A class of 30 students was surveyed: 18 play sport, 12 play an
  instrument, and 6 do both.

  \begin{enumerate}
  \def\labelenumii{\alph{enumii}.}
  \tightlist
  \item
    Draw a Venn diagram.
  \item
    Find P(plays sport but not instrument).
  \item
    Find P(neither sport nor instrument).
  \end{enumerate}
\item
  If P(A) = 0.4 and P(B) = 0.3 and A and B are independent, find P(A and
  B).
\end{enumerate}

\begin{tcolorbox}[enhanced jigsaw, bottomtitle=1mm, opacityback=0, titlerule=0mm, leftrule=.75mm, opacitybacktitle=0.6, coltitle=black, toprule=.15mm, colbacktitle=quarto-callout-note-color!10!white, title=\textcolor{quarto-callout-note-color}{\faInfo}\hspace{0.5em}{Answers}, left=2mm, rightrule=.15mm, toptitle=1mm, colframe=quarto-callout-note-color-frame, breakable, arc=.35mm, bottomrule=.15mm, colback=white]

\begin{enumerate}
\def\labelenumi{\arabic{enumi}.}
\tightlist
\item
  \begin{enumerate}
  \def\labelenumii{\alph{enumii}.}
  \tightlist
  \item
    4/10 = 0.4; b. 6/10 = 0.6; c.~6/10 = 0.6
  \end{enumerate}
\item
  \begin{enumerate}
  \def\labelenumii{\alph{enumii}.}
  \setcounter{enumii}{1}
  \tightlist
  \item
    12/30 = 0.4; c.~(30 − 18 − 12 + 6)/30 = 6/30 = 0.2
  \end{enumerate}
\item
  P(A and B) = 0.4 × 0.3 = 0.12 (for independent events)
\end{enumerate}

\end{tcolorbox}

\bookmarksetup{startatroot}

\chapter{Interpreting Statistical
Information}\label{interpreting-statistical-information}

\bookmarksetup{startatroot}

\chapter{Interpreting Statistical Information}\label{sec-inference}

A major focus of AS91946 is the ability to \textbf{critically read and
interpret} statistical information from a variety of sources ---
including newspapers, websites, government reports, and scientific
studies.

\section{Reading statistical claims}\label{reading-statistical-claims}

Statistical claims appear everywhere in everyday life. As a critical
consumer of information, you need to ask:

\begin{enumerate}
\def\labelenumi{\arabic{enumi}.}
\tightlist
\item
  \textbf{Who} collected the data and why?
\item
  \textbf{How} was the data collected (survey, experiment, census)?
\item
  \textbf{How many} people or things were measured (sample size)?
\item
  \textbf{Who} was sampled --- does the sample represent the population?
\item
  \textbf{What} do the numbers actually mean in context?
\end{enumerate}

\section{Sources of bias}\label{sources-of-bias}

\textbf{Bias} is a systematic error that causes results to be
unrepresentative.

\begin{longtable}[]{@{}
  >{\raggedright\arraybackslash}p{(\linewidth - 4\tabcolsep) * \real{0.3889}}
  >{\raggedright\arraybackslash}p{(\linewidth - 4\tabcolsep) * \real{0.3611}}
  >{\raggedright\arraybackslash}p{(\linewidth - 4\tabcolsep) * \real{0.2500}}@{}}
\toprule\noalign{}
\begin{minipage}[b]{\linewidth}\raggedright
Type of bias
\end{minipage} & \begin{minipage}[b]{\linewidth}\raggedright
Description
\end{minipage} & \begin{minipage}[b]{\linewidth}\raggedright
Example
\end{minipage} \\
\midrule\noalign{}
\endhead
\bottomrule\noalign{}
\endlastfoot
\textbf{Sampling bias} & Sample does not represent population &
Surveying only Year 11s about a school-wide issue \\
\textbf{Response bias} & People don't answer truthfully & ``Do you cheat
on tests?'' \\
\textbf{Non-response bias} & Certain groups are less likely to respond &
Online survey misses those without internet \\
\textbf{Leading question bias} & Question wording influences the answer
& ``Don't you agree that homework is too much?'' \\
\textbf{Volunteer bias} & Only motivated people respond & Self-selected
online polls \\
\end{longtable}

\begin{figure}

\centering{

\pandocbounded{\includegraphics[keepaspectratio]{chapters/06-inference_files/figure-latex/fig-bias-types-1.pdf}}

}

\caption{\label{fig-bias-types}Different types of bias and how they
affect results.}

\end{figure}%

\section{Sample size and reliability}\label{sample-size-and-reliability}

Larger samples generally give more \textbf{reliable} results --- but
size alone is not enough if the sample is biased.

\begin{Shaded}
\begin{Highlighting}[]
\FunctionTok{set.seed}\NormalTok{(}\DecValTok{42}\NormalTok{)}
\NormalTok{true\_prop }\OtherTok{\textless{}{-}} \FloatTok{0.55}  \CommentTok{\# True proportion in population}
\NormalTok{sample\_sizes }\OtherTok{\textless{}{-}} \FunctionTok{c}\NormalTok{(}\DecValTok{10}\NormalTok{, }\DecValTok{30}\NormalTok{, }\DecValTok{50}\NormalTok{, }\DecValTok{100}\NormalTok{, }\DecValTok{200}\NormalTok{, }\DecValTok{500}\NormalTok{, }\DecValTok{1000}\NormalTok{)}
\NormalTok{n\_simulations }\OtherTok{\textless{}{-}} \DecValTok{1000}

\FunctionTok{set.seed}\NormalTok{(}\DecValTok{321}\NormalTok{)}
\NormalTok{variability }\OtherTok{\textless{}{-}} \FunctionTok{sapply}\NormalTok{(sample\_sizes, }\ControlFlowTok{function}\NormalTok{(n) \{}
\NormalTok{  estimates }\OtherTok{\textless{}{-}} \FunctionTok{replicate}\NormalTok{(n\_simulations, }\FunctionTok{mean}\NormalTok{(}\FunctionTok{sample}\NormalTok{(}\FunctionTok{c}\NormalTok{(}\ConstantTok{TRUE}\NormalTok{, }\ConstantTok{FALSE}\NormalTok{), n,}
                                                    \AttributeTok{replace =} \ConstantTok{TRUE}\NormalTok{,}
                                                    \AttributeTok{prob =} \FunctionTok{c}\NormalTok{(true\_prop,}
                                                             \DecValTok{1} \SpecialCharTok{{-}}\NormalTok{ true\_prop))))}
  \FunctionTok{sd}\NormalTok{(estimates)}
\NormalTok{\})}

\FunctionTok{plot}\NormalTok{(}
\NormalTok{  sample\_sizes, variability,}
  \AttributeTok{type =} \StringTok{"b"}\NormalTok{,}
  \AttributeTok{pch  =} \DecValTok{19}\NormalTok{,}
  \AttributeTok{col  =} \StringTok{"\#0571b0"}\NormalTok{,}
  \AttributeTok{lwd  =} \DecValTok{2}\NormalTok{,}
  \AttributeTok{xlab =} \StringTok{"Sample Size (n)"}\NormalTok{,}
  \AttributeTok{ylab =} \StringTok{"Standard Deviation of Estimates"}\NormalTok{,}
  \AttributeTok{main =} \StringTok{"Sample Size vs Variability of Estimates"}\NormalTok{,}
  \AttributeTok{las  =} \DecValTok{1}
\NormalTok{)}
\FunctionTok{abline}\NormalTok{(}\AttributeTok{h =} \DecValTok{0}\NormalTok{, }\AttributeTok{col =} \StringTok{"\#d7191c"}\NormalTok{, }\AttributeTok{lty =} \DecValTok{2}\NormalTok{)}
\end{Highlighting}
\end{Shaded}

\begin{figure}[H]

\centering{

\pandocbounded{\includegraphics[keepaspectratio]{chapters/06-inference_files/figure-latex/fig-sample-size-1.pdf}}

}

\caption{\label{fig-sample-size}Variability in sample estimates
decreases as sample size increases.}

\end{figure}%

\section{Correlation and causation}\label{correlation-and-causation}

One of the most important concepts in statistical literacy is
understanding the difference between \textbf{correlation} and
\textbf{causation}.

\begin{tcolorbox}[enhanced jigsaw, bottomtitle=1mm, opacityback=0, titlerule=0mm, leftrule=.75mm, opacitybacktitle=0.6, coltitle=black, toprule=.15mm, colbacktitle=quarto-callout-warning-color!10!white, title=\textcolor{quarto-callout-warning-color}{\faExclamationTriangle}\hspace{0.5em}{Correlation ≠ Causation}, left=2mm, rightrule=.15mm, toptitle=1mm, colframe=quarto-callout-warning-color-frame, breakable, arc=.35mm, bottomrule=.15mm, colback=white]

Two variables can be \textbf{correlated} (move together) without one
causing the other. There may be a \textbf{confounding variable}
responsible for both.

\end{tcolorbox}

\subsection{Famous spurious
correlations}\label{famous-spurious-correlations}

\begin{itemize}
\tightlist
\item
  Ice cream sales and shark attacks both increase in summer → they are
  both caused by \textbf{hot weather}, not each other.
\item
  Countries with more TVs have longer life expectancy → both are driven
  by \textbf{wealth}, not TVs causing people to live longer.
\end{itemize}

\begin{figure}

\centering{

\pandocbounded{\includegraphics[keepaspectratio]{chapters/06-inference_files/figure-latex/fig-spurious-1.pdf}}

}

\caption{\label{fig-spurious}A spurious correlation: ice cream sales and
drownings both increase in summer.}

\end{figure}%

\section{Interpreting graphs
critically}\label{interpreting-graphs-critically}

When reading statistical graphs, watch out for:

\subsection{Misleading y-axis}\label{misleading-y-axis}

A y-axis that doesn't start at zero can make small differences look
huge.

\begin{Shaded}
\begin{Highlighting}[]
\FunctionTok{par}\NormalTok{(}\AttributeTok{mfrow =} \FunctionTok{c}\NormalTok{(}\DecValTok{1}\NormalTok{, }\DecValTok{2}\NormalTok{))}

\NormalTok{years   }\OtherTok{\textless{}{-}} \DecValTok{2020}\SpecialCharTok{:}\DecValTok{2024}
\NormalTok{revenue }\OtherTok{\textless{}{-}} \FunctionTok{c}\NormalTok{(}\DecValTok{980}\NormalTok{, }\DecValTok{990}\NormalTok{, }\DecValTok{995}\NormalTok{, }\DecValTok{1005}\NormalTok{, }\DecValTok{1010}\NormalTok{)}

\CommentTok{\# Misleading}
\FunctionTok{plot}\NormalTok{(years, revenue, }\AttributeTok{type =} \StringTok{"b"}\NormalTok{, }\AttributeTok{pch =} \DecValTok{19}\NormalTok{, }\AttributeTok{col =} \StringTok{"\#d7191c"}\NormalTok{, }\AttributeTok{lwd =} \DecValTok{2}\NormalTok{,}
     \AttributeTok{ylim =} \FunctionTok{c}\NormalTok{(}\DecValTok{975}\NormalTok{, }\DecValTok{1015}\NormalTok{), }\AttributeTok{xlab =} \StringTok{"Year"}\NormalTok{, }\AttributeTok{ylab =} \StringTok{"Revenue ($M)"}\NormalTok{,}
     \AttributeTok{main =} \StringTok{"Misleading: Truncated Axis"}\NormalTok{, }\AttributeTok{las =} \DecValTok{1}\NormalTok{)}

\CommentTok{\# Honest}
\FunctionTok{plot}\NormalTok{(years, revenue, }\AttributeTok{type =} \StringTok{"b"}\NormalTok{, }\AttributeTok{pch =} \DecValTok{19}\NormalTok{, }\AttributeTok{col =} \StringTok{"\#4dac26"}\NormalTok{, }\AttributeTok{lwd =} \DecValTok{2}\NormalTok{,}
     \AttributeTok{ylim =} \FunctionTok{c}\NormalTok{(}\DecValTok{0}\NormalTok{, }\DecValTok{1100}\NormalTok{), }\AttributeTok{xlab =} \StringTok{"Year"}\NormalTok{, }\AttributeTok{ylab =} \StringTok{"Revenue ($M)"}\NormalTok{,}
     \AttributeTok{main =} \StringTok{"Honest: Starting at Zero"}\NormalTok{, }\AttributeTok{las =} \DecValTok{1}\NormalTok{)}

\FunctionTok{par}\NormalTok{(}\AttributeTok{mfrow =} \FunctionTok{c}\NormalTok{(}\DecValTok{1}\NormalTok{, }\DecValTok{1}\NormalTok{))}
\end{Highlighting}
\end{Shaded}

\begin{figure}[H]

\centering{

\pandocbounded{\includegraphics[keepaspectratio]{chapters/06-inference_files/figure-latex/fig-misleading-axis-1.pdf}}

}

\caption{\label{fig-misleading-axis}The same data displayed with a
truncated y-axis (left) vs starting at zero (right).}

\end{figure}%

\subsection{Percentage vs absolute
numbers}\label{percentage-vs-absolute-numbers}

\begin{tcolorbox}[enhanced jigsaw, bottomtitle=1mm, opacityback=0, titlerule=0mm, leftrule=.75mm, opacitybacktitle=0.6, coltitle=black, toprule=.15mm, colbacktitle=quarto-callout-tip-color!10!white, title=\textcolor{quarto-callout-tip-color}{\faLightbulb}\hspace{0.5em}{Be careful with percentages}, left=2mm, rightrule=.15mm, toptitle=1mm, colframe=quarto-callout-tip-color-frame, breakable, arc=.35mm, bottomrule=.15mm, colback=white]

``Crime increased by 100\%'' sounds alarming --- but if there were only
2 crimes last year and 4 this year, the absolute increase is tiny.

Always ask: \textbf{100\% of what?}

\end{tcolorbox}

\subsection{Averages without context}\label{averages-without-context}

The \textbf{mean} can be pulled by extreme values, making it misleading
if used alone. Always ask what \textbf{measure of centre} was used and
whether it is appropriate.

\section{Reading media reports --- a
framework}\label{reading-media-reports-a-framework}

Use these questions when evaluating a statistical claim in the media:

\begin{enumerate}
\def\labelenumi{\arabic{enumi}.}
\tightlist
\item
  \textbf{What is being measured?} (the variable and its units)
\item
  \textbf{Who was studied?} (the population and sample)
\item
  \textbf{How was data collected?} (method and potential biases)
\item
  \textbf{How large was the sample?}
\item
  \textbf{What is the source?} (government, independent researcher,
  company with a vested interest?)
\item
  \textbf{Are conclusions supported by the evidence?} Or is the claim
  overstated?
\item
  \textbf{Are comparisons fair?} Are groups being compared actually
  comparable?
\end{enumerate}

\section{A worked example}\label{a-worked-example}

\begin{tcolorbox}[enhanced jigsaw, bottomtitle=1mm, opacityback=0, titlerule=0mm, leftrule=.75mm, opacitybacktitle=0.6, coltitle=black, toprule=.15mm, colbacktitle=quarto-callout-note-color!10!white, title=\textcolor{quarto-callout-note-color}{\faInfo}\hspace{0.5em}{News headline}, left=2mm, rightrule=.15mm, toptitle=1mm, colframe=quarto-callout-note-color-frame, breakable, arc=.35mm, bottomrule=.15mm, colback=white]

\emph{``Study shows eating chocolate improves exam scores by 15\%''}

\end{tcolorbox}

\textbf{Critical analysis:}

\begin{itemize}
\tightlist
\item
  \textbf{What was measured?} Exam scores before and after eating
  chocolate.
\item
  \textbf{Who was studied?} Not stated --- red flag!
\item
  \textbf{Sample size?} Not given --- could be only 10 students.
\item
  \textbf{Confounding variables?} Students who eat chocolate may be from
  wealthier families with better study resources.
\item
  \textbf{Causation or correlation?} This is a correlational study;
  chocolate doesn't necessarily \emph{cause} better scores.
\item
  \textbf{Who funded it?} If a chocolate company funded it, there is a
  conflict of interest.
\end{itemize}

\textbf{Conclusion:} The headline is likely exaggerated or misleading.
We would need much more information before accepting this claim.

\section{Glossary}\label{glossary-4}

\begin{longtable}[]{@{}
  >{\raggedright\arraybackslash}p{(\linewidth - 2\tabcolsep) * \real{0.3333}}
  >{\raggedright\arraybackslash}p{(\linewidth - 2\tabcolsep) * \real{0.6667}}@{}}
\toprule\noalign{}
\begin{minipage}[b]{\linewidth}\raggedright
Term
\end{minipage} & \begin{minipage}[b]{\linewidth}\raggedright
Definition
\end{minipage} \\
\midrule\noalign{}
\endhead
\bottomrule\noalign{}
\endlastfoot
\textbf{Bias} & Systematic error that distorts results \\
\textbf{Confounding variable} & A hidden variable that explains an
apparent relationship \\
\textbf{Correlation} & A statistical relationship between two
variables \\
\textbf{Causation} & One variable directly causes a change in another \\
\textbf{Reliable} & Consistent results if repeated \\
\textbf{Valid} & Measuring what it claims to measure \\
\textbf{Population} & The full group of interest \\
\textbf{Sample} & A subset of the population \\
\end{longtable}

\section{Check your understanding}\label{check-your-understanding-5}

\begin{enumerate}
\def\labelenumi{\arabic{enumi}.}
\item
  A newspaper article claims: ``People who eat breakfast earn more money
  than those who skip it.'' Identify a possible confounding variable.
\item
  A survey found that 80\% of students support a longer lunch break. The
  survey was conducted by standing outside the school canteen. What type
  of bias is this? How might it affect results?
\item
  A graph shows the number of electric vehicles sold in New Zealand,
  with the y-axis starting at 45,000 rather than 0. Sales increased from
  47,000 to 51,000. On the graph, this looks like a 400\% increase.
  Calculate the true percentage increase.
\item
  Explain the difference between correlation and causation using your
  own example.
\end{enumerate}

\begin{tcolorbox}[enhanced jigsaw, bottomtitle=1mm, opacityback=0, titlerule=0mm, leftrule=.75mm, opacitybacktitle=0.6, coltitle=black, toprule=.15mm, colbacktitle=quarto-callout-note-color!10!white, title=\textcolor{quarto-callout-note-color}{\faInfo}\hspace{0.5em}{Answers}, left=2mm, rightrule=.15mm, toptitle=1mm, colframe=quarto-callout-note-color-frame, breakable, arc=.35mm, bottomrule=.15mm, colback=white]

\begin{enumerate}
\def\labelenumi{\arabic{enumi}.}
\tightlist
\item
  Wealthier people are more likely to eat breakfast (they can afford
  food) and also earn more --- wealth is the confounding variable.
\item
  Sampling bias / volunteer bias --- students near the canteen are more
  likely to want a longer lunch break. Results will overestimate
  support.
\item
  True \% increase = (51,000 − 47,000) / 47,000 × 100 ≈ 8.5\%.
\item
  Accept any reasonable example, e.g., ``Countries with more hospitals
  have higher death rates --- hospitals don't cause death; both are
  linked to population size.''
\end{enumerate}

\end{tcolorbox}

\bookmarksetup{startatroot}

\chapter{Practice Questions and Exam
Tips}\label{practice-questions-and-exam-tips}

\bookmarksetup{startatroot}

\chapter{Practice Questions and Achievement Standard
Tips}\label{sec-practice}

This chapter brings everything together with practice questions at each
grade level, and advice on how to approach the AS91946 assessment.

\section{About the achievement
standard}\label{about-the-achievement-standard}

\begin{longtable}[]{@{}
  >{\raggedright\arraybackslash}p{(\linewidth - 2\tabcolsep) * \real{0.3810}}
  >{\raggedright\arraybackslash}p{(\linewidth - 2\tabcolsep) * \real{0.6190}}@{}}
\toprule\noalign{}
\begin{minipage}[b]{\linewidth}\raggedright
Detail
\end{minipage} & \begin{minipage}[b]{\linewidth}\raggedright
Information
\end{minipage} \\
\midrule\noalign{}
\endhead
\bottomrule\noalign{}
\endlastfoot
\textbf{Standard number} & AS91946 \\
\textbf{Title} & Interpret and apply mathematical and statistical
information in context \\
\textbf{Level} & NCEA Level 1 \\
\textbf{Credits} & 4 \\
\textbf{Assessment} & External (end-of-year exam) \\
\textbf{Strands} & Statistics and Probability \\
\end{longtable}

\section{What each grade requires}\label{what-each-grade-requires}

\subsection{Achieved}\label{achieved}

\begin{itemize}
\tightlist
\item
  Correctly read values from graphs and tables
\item
  Calculate basic statistics (mean, median, mode, range)
\item
  State probabilities from given information
\item
  Identify variable types
\item
  Make simple statements about data in context
\end{itemize}

\subsection{Merit}\label{merit}

\begin{itemize}
\tightlist
\item
  \textbf{Justify} your answers with reasons
\item
  Compare and contrast groups using both centre AND spread
\item
  Identify potential sources of bias
\item
  Explain the difference between correlation and causation
\item
  Interpret probabilities and relate them to real-world context
\end{itemize}

\subsection{Excellence}\label{excellence}

\begin{itemize}
\tightlist
\item
  Show \textbf{statistical insight} --- go beyond the obvious
\item
  Consider alternative explanations or limitations
\item
  Make connections between different parts of the data
\item
  Critically evaluate statistical claims
\item
  Use appropriate statistical language precisely
\end{itemize}

\section{Writing tips for the exam}\label{writing-tips-for-the-exam}

\begin{tcolorbox}[enhanced jigsaw, bottomtitle=1mm, opacityback=0, titlerule=0mm, leftrule=.75mm, opacitybacktitle=0.6, coltitle=black, toprule=.15mm, colbacktitle=quarto-callout-tip-color!10!white, title=\textcolor{quarto-callout-tip-color}{\faLightbulb}\hspace{0.5em}{The golden rule}, left=2mm, rightrule=.15mm, toptitle=1mm, colframe=quarto-callout-tip-color-frame, breakable, arc=.35mm, bottomrule=.15mm, colback=white]

\textbf{Always refer back to the context.} Every answer should mention:
- The \textbf{variable} being described - The \textbf{units} of
measurement - The \textbf{group} being described (population/sample)

❌ ``The mean is 65.''\\
✅ ``The mean maths score for Year 11 students was 65 out of 100.''

\end{tcolorbox}

\subsection{Comparison sentence
starters}\label{comparison-sentence-starters}

\begin{longtable}[]{@{}
  >{\raggedright\arraybackslash}p{(\linewidth - 2\tabcolsep) * \real{0.3333}}
  >{\raggedright\arraybackslash}p{(\linewidth - 2\tabcolsep) * \real{0.6667}}@{}}
\toprule\noalign{}
\begin{minipage}[b]{\linewidth}\raggedright
Purpose
\end{minipage} & \begin{minipage}[b]{\linewidth}\raggedright
Sentence starters
\end{minipage} \\
\midrule\noalign{}
\endhead
\bottomrule\noalign{}
\endlastfoot
Centre & ``On average\ldots{}''; ``The typical\ldots{}''; ``The
median\ldots is higher/lower than\ldots{}'' \\
Spread & ``The scores are more/less variable\ldots{}''; ``The IQR is
larger, suggesting\ldots{}'' \\
Shape & ``The distribution is roughly symmetric /
right-skewed\ldots{}'' \\
Outliers & ``There appears to be an outlier at\ldots which may
indicate\ldots{}'' \\
Causation & ``However, this may be due to\ldots rather than a direct
causal relationship'' \\
\end{longtable}

\section{Practice set 1 --- Achieved
level}\label{practice-set-1-achieved-level}

\begin{verbatim}
#> Heights (cm) of 20 Year 11 students:
#> 147.3 153.9 154.5 157.3 158.9 159.5 163.6 164.3 164.8 165 165.2 166.4 166.5 167.2 169.3 170.1 171.3 171.5 175.2 184.6
\end{verbatim}

\textbf{Questions:}

\begin{enumerate}
\def\labelenumi{\arabic{enumi}.}
\tightlist
\item
  What type of data is height? Give a reason.
\item
  Calculate the mean height. Round to one decimal place.
\item
  Find the median height.
\item
  Find the range.
\item
  A student drew a histogram of this data. Circle the most appropriate
  description: \textbf{symmetric / right-skewed / left-skewed}
\end{enumerate}

\begin{tcolorbox}[enhanced jigsaw, bottomtitle=1mm, opacityback=0, titlerule=0mm, leftrule=.75mm, opacitybacktitle=0.6, coltitle=black, toprule=.15mm, colbacktitle=quarto-callout-note-color!10!white, title=\textcolor{quarto-callout-note-color}{\faInfo}\hspace{0.5em}{Worked solutions}, left=2mm, rightrule=.15mm, toptitle=1mm, colframe=quarto-callout-note-color-frame, breakable, arc=.35mm, bottomrule=.15mm, colback=white]

\begin{Shaded}
\begin{Highlighting}[]
\FunctionTok{set.seed}\NormalTok{(}\DecValTok{111}\NormalTok{)}
\NormalTok{heights }\OtherTok{\textless{}{-}} \FunctionTok{round}\NormalTok{(}\FunctionTok{rnorm}\NormalTok{(}\DecValTok{20}\NormalTok{, }\DecValTok{168}\NormalTok{, }\DecValTok{9}\NormalTok{), }\DecValTok{1}\NormalTok{)}

\FunctionTok{cat}\NormalTok{(}\StringTok{"1. Numerical, continuous — height can take any value within a range.}\SpecialCharTok{\textbackslash{}n\textbackslash{}n}\StringTok{"}\NormalTok{)}
\CommentTok{\#\textgreater{} 1. Numerical, continuous — height can take any value within a range.}
\FunctionTok{cat}\NormalTok{(}\StringTok{"2. Mean:"}\NormalTok{, }\FunctionTok{round}\NormalTok{(}\FunctionTok{mean}\NormalTok{(heights), }\DecValTok{1}\NormalTok{), }\StringTok{"cm}\SpecialCharTok{\textbackslash{}n\textbackslash{}n}\StringTok{"}\NormalTok{)}
\CommentTok{\#\textgreater{} 2. Mean: 164.8 cm}
\FunctionTok{cat}\NormalTok{(}\StringTok{"3. Median:"}\NormalTok{, }\FunctionTok{median}\NormalTok{(heights), }\StringTok{"cm}\SpecialCharTok{\textbackslash{}n\textbackslash{}n}\StringTok{"}\NormalTok{)}
\CommentTok{\#\textgreater{} 3. Median: 165.1 cm}
\FunctionTok{cat}\NormalTok{(}\StringTok{"4. Range:"}\NormalTok{, }\FunctionTok{max}\NormalTok{(heights) }\SpecialCharTok{{-}} \FunctionTok{min}\NormalTok{(heights), }\StringTok{"cm}\SpecialCharTok{\textbackslash{}n\textbackslash{}n}\StringTok{"}\NormalTok{)}
\CommentTok{\#\textgreater{} 4. Range: 37.3 cm}
\FunctionTok{cat}\NormalTok{(}\StringTok{"5. Likely symmetric (heights tend to be normally distributed).}\SpecialCharTok{\textbackslash{}n}\StringTok{"}\NormalTok{)}
\CommentTok{\#\textgreater{} 5. Likely symmetric (heights tend to be normally distributed).}
\end{Highlighting}
\end{Shaded}

\end{tcolorbox}

\section{Practice set 2 --- Merit
level}\label{practice-set-2-merit-level}

A local council surveyed 80 residents about their weekly exercise hours.
The results were:

\begin{Shaded}
\begin{Highlighting}[]
\FunctionTok{set.seed}\NormalTok{(}\DecValTok{222}\NormalTok{)}
\NormalTok{adults    }\OtherTok{\textless{}{-}} \FunctionTok{round}\NormalTok{(}\FunctionTok{pmax}\NormalTok{(}\DecValTok{0}\NormalTok{, }\FunctionTok{rnorm}\NormalTok{(}\DecValTok{40}\NormalTok{, }\AttributeTok{mean =} \FloatTok{4.5}\NormalTok{, }\AttributeTok{sd =} \FloatTok{2.2}\NormalTok{)), }\DecValTok{1}\NormalTok{)}
\NormalTok{teenagers }\OtherTok{\textless{}{-}} \FunctionTok{round}\NormalTok{(}\FunctionTok{pmax}\NormalTok{(}\DecValTok{0}\NormalTok{, }\FunctionTok{rnorm}\NormalTok{(}\DecValTok{40}\NormalTok{, }\AttributeTok{mean =} \FloatTok{6.2}\NormalTok{, }\AttributeTok{sd =} \FloatTok{3.1}\NormalTok{)), }\DecValTok{1}\NormalTok{)}

\NormalTok{summary\_tbl }\OtherTok{\textless{}{-}} \FunctionTok{data.frame}\NormalTok{(}
  \AttributeTok{Group      =} \FunctionTok{c}\NormalTok{(}\StringTok{"Adults"}\NormalTok{, }\StringTok{"Teenagers"}\NormalTok{),}
  \AttributeTok{n          =} \DecValTok{40}\NormalTok{,}
  \AttributeTok{Mean       =} \FunctionTok{round}\NormalTok{(}\FunctionTok{c}\NormalTok{(}\FunctionTok{mean}\NormalTok{(adults), }\FunctionTok{mean}\NormalTok{(teenagers)), }\DecValTok{2}\NormalTok{),}
  \AttributeTok{Median     =} \FunctionTok{c}\NormalTok{(}\FunctionTok{median}\NormalTok{(adults), }\FunctionTok{median}\NormalTok{(teenagers)),}
  \AttributeTok{SD         =} \FunctionTok{round}\NormalTok{(}\FunctionTok{c}\NormalTok{(}\FunctionTok{sd}\NormalTok{(adults), }\FunctionTok{sd}\NormalTok{(teenagers)), }\DecValTok{2}\NormalTok{),}
  \AttributeTok{IQR        =} \FunctionTok{round}\NormalTok{(}\FunctionTok{c}\NormalTok{(}\FunctionTok{IQR}\NormalTok{(adults), }\FunctionTok{IQR}\NormalTok{(teenagers)), }\DecValTok{2}\NormalTok{),}
  \AttributeTok{Min        =} \FunctionTok{c}\NormalTok{(}\FunctionTok{min}\NormalTok{(adults), }\FunctionTok{min}\NormalTok{(teenagers)),}
  \AttributeTok{Max        =} \FunctionTok{c}\NormalTok{(}\FunctionTok{max}\NormalTok{(adults), }\FunctionTok{max}\NormalTok{(teenagers))}
\NormalTok{)}

\NormalTok{knitr}\SpecialCharTok{::}\FunctionTok{kable}\NormalTok{(summary\_tbl, }\AttributeTok{caption =} \StringTok{"Weekly exercise hours: adults vs teenagers"}\NormalTok{)}
\end{Highlighting}
\end{Shaded}

\begin{longtable}[]{@{}lrrrrrrr@{}}
\caption{Weekly exercise hours: adults vs teenagers}\tabularnewline
\toprule\noalign{}
Group & n & Mean & Median & SD & IQR & Min & Max \\
\midrule\noalign{}
\endfirsthead
\toprule\noalign{}
Group & n & Mean & Median & SD & IQR & Min & Max \\
\midrule\noalign{}
\endhead
\bottomrule\noalign{}
\endlastfoot
Adults & 40 & 4.45 & 4.35 & 2.07 & 2.85 & 0 & 7.9 \\
Teenagers & 40 & 6.58 & 7.10 & 2.85 & 3.58 & 0 & 13.0 \\
\end{longtable}

\begin{Shaded}
\begin{Highlighting}[]
\FunctionTok{boxplot}\NormalTok{(}
  \FunctionTok{list}\NormalTok{(}\AttributeTok{Adults =}\NormalTok{ adults, }\AttributeTok{Teenagers =}\NormalTok{ teenagers),}
  \AttributeTok{col    =} \FunctionTok{c}\NormalTok{(}\StringTok{"\#4dac26"}\NormalTok{, }\StringTok{"\#0571b0"}\NormalTok{),}
   \AttributeTok{border =} \StringTok{"\#333333"}\NormalTok{,}
  \AttributeTok{main   =} \StringTok{"Weekly Exercise Hours"}\NormalTok{,}
  \AttributeTok{ylab   =} \StringTok{"Hours per week"}\NormalTok{,}
  \AttributeTok{las    =} \DecValTok{1}
\NormalTok{)}
\end{Highlighting}
\end{Shaded}

\begin{figure}[H]

\centering{

\pandocbounded{\includegraphics[keepaspectratio]{chapters/07-practice_files/figure-latex/fig-merit-boxplot-1.pdf}}

}

\caption{\label{fig-merit-boxplot}Weekly exercise hours for adults and
teenagers.}

\end{figure}%

\textbf{Questions:}

\begin{enumerate}
\def\labelenumi{\arabic{enumi}.}
\tightlist
\item
  Compare the centre of the two distributions. What does this tell us
  about exercise habits in context?
\item
  Compare the spread. What does this tell us?
\item
  The council used a telephone survey conducted at 2 pm on weekdays.
  Identify one source of bias and explain how it might affect the
  results.
\item
  Can we conclude from this data that teenagers exercise more because
  they are younger? Explain.
\end{enumerate}

\begin{tcolorbox}[enhanced jigsaw, bottomtitle=1mm, opacityback=0, titlerule=0mm, leftrule=.75mm, opacitybacktitle=0.6, coltitle=black, toprule=.15mm, colbacktitle=quarto-callout-note-color!10!white, title=\textcolor{quarto-callout-note-color}{\faInfo}\hspace{0.5em}{Model answers}, left=2mm, rightrule=.15mm, toptitle=1mm, colframe=quarto-callout-note-color-frame, breakable, arc=.35mm, bottomrule=.15mm, colback=white]

\textbf{1. Centre:} The median exercise hours for teenagers (approx. 7.1
hrs) is higher than for adults (4.35 hrs). This suggests teenagers in
this sample tend to do more weekly exercise than adults.

\textbf{2. Spread:} The IQR and standard deviation are larger for
teenagers (3.6 hrs vs 2.9 hrs). This means teenagers' exercise habits
are more variable --- some exercise a lot while others exercise very
little. Adults' exercise habits are more consistent.

\textbf{3. Bias:} Calling at 2 pm on weekdays means working adults are
unlikely to answer (they are at work). This leads to
\textbf{non-response bias} --- the adult sample may over-represent
retired or part-time workers, who may exercise more than the general
adult population.

\textbf{4. Causation:} No --- we cannot conclude that age \emph{causes}
the difference in exercise. Teenagers may exercise more because they
participate in school sport, not because of their age per se. Age and
exercise are \emph{correlated} in this sample, but there may be
confounding variables such as school PE requirements, free time, or
lifestyle factors.

\end{tcolorbox}

\section{Practice set 3 --- Excellence
level}\label{practice-set-3-excellence-level}

A researcher claims: \emph{``A study of 50 Auckland households found
that households with more bookshelves have higher household income.
Therefore, buying more bookshelves will increase your income.''}

\textbf{Questions:}

\begin{enumerate}
\def\labelenumi{\arabic{enumi}.}
\tightlist
\item
  Identify the two variables and their types.
\item
  Explain why this conclusion is flawed. Use statistical language.
\item
  Suggest one confounding variable that could explain the relationship.
\item
  What would the researcher need to do to properly establish causation?
\item
  The sample was 50 Auckland households selected by going door-to-door
  in Remuera (a wealthy suburb). Evaluate the reliability of
  generalising these findings to all New Zealand households.
\end{enumerate}

\begin{tcolorbox}[enhanced jigsaw, bottomtitle=1mm, opacityback=0, titlerule=0mm, leftrule=.75mm, opacitybacktitle=0.6, coltitle=black, toprule=.15mm, colbacktitle=quarto-callout-note-color!10!white, title=\textcolor{quarto-callout-note-color}{\faInfo}\hspace{0.5em}{Model answers}, left=2mm, rightrule=.15mm, toptitle=1mm, colframe=quarto-callout-note-color-frame, breakable, arc=.35mm, bottomrule=.15mm, colback=white]

\textbf{1.} Number of bookshelves: numerical, discrete. Household
income: numerical, continuous.

\textbf{2.} The conclusion confuses \textbf{correlation with causation}.
Even if the two variables are positively correlated, this does not mean
bookshelves cause higher income. The relationship could be due to a
third variable, and correlation studies cannot establish
cause-and-effect without experimental control.

\textbf{3.} \textbf{Wealth/education level} --- wealthier, more educated
households are more likely to own books/bookshelves AND to have higher
income. Both variables are driven by a common factor.

\textbf{4.} To establish causation, the researcher would need a
\textbf{controlled experiment} --- randomly assigning households to
receive bookshelves or not, holding all other variables constant, and
measuring income change over time. This would be impractical.

\textbf{5.} The sample has significant \textbf{sampling bias}. Remuera
is one of Auckland's wealthiest suburbs. The sample does not represent
the diversity of New Zealand households in terms of income, region, or
ethnicity. Results cannot be reliably generalised to all New Zealand
households. A \textbf{stratified random sample} from across New Zealand
would be needed.

\end{tcolorbox}

\section{Key formulas summary}\label{key-formulas-summary}

\begin{longtable}[]{@{}ll@{}}
\toprule\noalign{}
Measure & Formula \\
\midrule\noalign{}
\endhead
\bottomrule\noalign{}
\endlastfoot
Mean & x̄ = Σxᵢ / n \\
Range & Max − Min \\
IQR & Q3 − Q1 \\
Probability & P(A) = favourable outcomes / total outcomes \\
Complement & P(A') = 1 − P(A) \\
Lower outlier fence & Q1 − 1.5 × IQR \\
Upper outlier fence & Q3 + 1.5 × IQR \\
\end{longtable}

\section{Exam strategy checklist}\label{exam-strategy-checklist}

Before the exam, make sure you can:

\begin{itemize}
\tightlist
\item[$\square$]
  Identify variable types (categorical/numerical, nominal/ordinal,
  discrete/continuous)
\item[$\square$]
  Read values accurately from bar charts, histograms, box plots, and
  scatter plots
\item[$\square$]
  Calculate mean, median, mode, range, IQR
\item[$\square$]
  Describe distributions using SOCS (Shape, Outliers, Centre, Spread)
\item[$\square$]
  Compare two groups using both centre and spread in context
\item[$\square$]
  Calculate simple probabilities and use complementary events
\item[$\square$]
  Read and calculate from two-way tables
\item[$\square$]
  Identify bias in data collection methods
\item[$\square$]
  Distinguish correlation from causation
\item[$\square$]
  Critically evaluate statistical claims in media
\end{itemize}

\section{Useful resources}\label{useful-resources}

\begin{longtable}[]{@{}
  >{\raggedright\arraybackslash}p{(\linewidth - 4\tabcolsep) * \real{0.3571}}
  >{\raggedright\arraybackslash}p{(\linewidth - 4\tabcolsep) * \real{0.4643}}
  >{\raggedright\arraybackslash}p{(\linewidth - 4\tabcolsep) * \real{0.1786}}@{}}
\toprule\noalign{}
\begin{minipage}[b]{\linewidth}\raggedright
Resource
\end{minipage} & \begin{minipage}[b]{\linewidth}\raggedright
Description
\end{minipage} & \begin{minipage}[b]{\linewidth}\raggedright
URL
\end{minipage} \\
\midrule\noalign{}
\endhead
\bottomrule\noalign{}
\endlastfoot
NZQA & Official achievement standard information &
\href{https://www.nzqa.govt.nz}{nzqa.govt.nz} \\
CensusAtSchool & Real NZ student data for practice &
\href{https://new.censusatschool.org.nz}{new.censusatschool.org.nz} \\
iNZight & Free NZ-made statistics software &
\href{https://inzight.nz}{inzight.nz} \\
Stats NZ & Official NZ statistics and data &
\href{https://www.stats.govt.nz}{stats.govt.nz} \\
Khan Academy & Free video lessons on statistics &
\href{https://www.khanacademy.org}{khanacademy.org} \\
\end{longtable}

\bookmarksetup{startatroot}

\chapter{References}\label{references}

\bookmarksetup{startatroot}

\chapter*{References}\label{references-1}
\addcontentsline{toc}{chapter}{References}

\markboth{References}{References}

\section{Achievement Standard}\label{achievement-standard}

New Zealand Qualifications Authority (NZQA). (2024). \emph{Achievement
Standard 91946: Interpret and apply mathematical and statistical
information in context.} NZQA.
\url{https://www.nzqa.govt.nz/nqfdocs/ncea-resource/achievements/2024/as91946.pdf}

\section{New Zealand Statistics
Curriculum}\label{new-zealand-statistics-curriculum}

Ministry of Education New Zealand. (2007). \emph{The New Zealand
Curriculum.} Learning Media. \url{https://nzcurriculum.tki.org.nz}

Ministry of Education New Zealand. (2023). \emph{Mathematics and
Statistics --- Achievement Objectives.}
\url{https://nzcurriculum.tki.org.nz/The-New-Zealand-Curriculum/Mathematics-and-statistics}

\section{Textbooks and Teaching
Resources}\label{textbooks-and-teaching-resources}

Wild, C. J., \& Seber, G. A. F. (2000). \emph{Chance Encounters: A First
Course in Data Analysis and Inference.} Wiley.

Pfannkuch, M., \& Wild, C. J. (2004). Towards an understanding of
statistical thinking. In D. Ben-Zvi \& J. Garfield (Eds.), \emph{The
Challenge of Developing Statistical Literacy, Reasoning and Thinking}
(pp.~17--46). Kluwer Academic.

\section{Data Sources}\label{data-sources}

Statistics New Zealand (Stats NZ). (2024). \emph{Household Economic
Survey.} \url{https://www.stats.govt.nz}

CensusAtSchool New Zealand. (2024). \emph{CensusAtSchool data.}
University of Auckland. \url{https://new.censusatschool.org.nz}

\section{Software}\label{software}

R Core Team. (2024). \emph{R: A Language and Environment for Statistical
Computing.} R Foundation for Statistical Computing.
\url{https://www.r-project.org}

Allaire, J. J., Teague, C., Scheidegger, C., Xie, Y., \& Dervieux, C.
(2024). \emph{Quarto} (Version 1.5). \url{https://quarto.org}

Wickham, H., et al.~(2023). \emph{ggplot2: Elegant Graphics for Data
Analysis} (3rd ed.). Springer. \url{https://ggplot2-book.org}

\section{Statistical Literacy}\label{statistical-literacy}

Huff, D. (1954). \emph{How to Lie with Statistics.} W. W. Norton \&
Company.

Wheelan, C. (2013). \emph{Naked Statistics: Stripping the Dread from the
Data.} W. W. Norton \& Company.


\backmatter


\end{document}
